% !TEX root = ../main.tex
% Chapter 1: Background.
% Included from main.tex via % !TEX root = ../main.tex
% Chapter 1: Background.
% Included from main.tex via % !TEX root = ../main.tex
% Chapter 1: Background.
% Included from main.tex via % !TEX root = ../main.tex
% Chapter 1: Background.
% Included from main.tex via \input{chapters/Chapter1}.

\chapter{Background}
\label{sec:background}
Before we talk about the main object in this thesis, we need to first introduce some analysis
backgrounds.
\section{Poisson integral}
Let $\mathbb{D}:=\{z\in\C: |z|<1\}$ and $\mathbb{T}=\partial \mathbb{D}$. 
When given a continuous function $f$ on $\mathbb{T}$, the famous \textbf{Dirichlet
problem}, which is to find the solution for the following equations:
\begin{equation*}
  \begin{cases}
  \Delta u=0, & z\in \mathbb{D}\\
  u=f, & z\in \mathbb{T}.
  \end{cases}
\end{equation*}
is asking for a harmonic extension $u$ over $\mathbb{D}$ with the boundary data $f$ on $\mathbb{T}$.
To give a solution for this problem, we'd like to introduce the notion called \textbf{Poisson integral}.
This section is based on \cite{garnett2006bounded}, I.3.
\begin{definition}
Let $c\in\mathbb{C}$. The \textbf{Herglotz--Riesz kernel centered at $c$} is the function
\begin{equation*}
  K(c;z,w) := \frac{(w-c)+(z-c)}{(w-c)-(z-c)} = \frac{w+z-2c}{w-z},
  \qquad z,w\in\mathbb{C},\ w\neq z.
\end{equation*}
The \textbf{Poisson kernel centered at $c$} is defined as
\begin{equation*}
  P(c;z,w) := \text{Re}\; K(c;z,w)= \frac{|w-c|^2-|z-c|^2}{|w-z|^2},
  \qquad z,w\in\mathbb{C},\ w\neq z.
\end{equation*}
and we denote $K(z,w):= K(0;z,w)$, $P(z,w):= P(0;z,w)$, and $K_z(\theta):= K(z,e^{i\theta})$,
 $P_z(\theta):= P(z,e^{i\theta})$. 
\end{definition}

\begin{definition}
Given a function $f(\theta)\in L^1(\mathbb{T})$, we define the \textbf{Poisson integral} of $f$ 
to be the function
\begin{equation*}
  P[f](z):=\frac{1}{2\pi}\int P_z(\theta) f(\theta)d\theta.
\end{equation*}
In general, given a finite (signed) measure $\mu$ on $\mathbb{T}$, we define the \textbf{Poisson integral} of $\mu$
to be the function
\begin{equation*}
  P[\mu](z):=\int P_z(\theta) d\mu (\theta).
\end{equation*}
The first definition is a special case of the second one if we set 
$d\mu(\theta)=\frac{1}{2\pi} f(\theta)d\theta$.
\end{definition}

\begin{theorem}[Poisson integral formula]\label{Poisson integral formula}
Let $u$ be a continuous function on the closed unit disc $\overline{\mathbb{D}}$. If $u$ is
harmonic on $\mathbb{D}$, then
\begin{equation*}
  u(z)=\frac{1}{2\pi}\int_0^{2\pi} P_z(\theta)\, u(e^{i\theta})\, d\theta
  \qquad \text{for all } z\in\mathbb{D}.
\end{equation*}
\end{theorem}
\begin{proof}
Recall that if $u$ is harmonic on $\mathbb{D}$, then it has the 
\textbf{mean value property}: 
\begin{equation*}
  u(0)=\frac{1}{2\pi}\int_0^{2\pi} u(e^{i\theta})d\theta
\end{equation*}
Now given $z_0=re^{i\theta_0}$, consider the M\"{o}bius transformation 
$\tau(z)=(z-z_0)/(1-\bar{z}_0 z)$, and write $\tau(e^{i\theta})=e^{i\varphi}$.
Then we can compute that $d\varphi/d\theta=P_{z_0}(\theta)$, therefore
\begin{equation*}
  u(z_0)=u(\tau^{-1}(0))=\frac{1}{2\pi}\int_0^{2\pi} (u\circ\tau^{-1})(e^{i\varphi})
  d\varphi=\frac{1}{2\pi}\int_0^{2\pi} u(e^{i\theta})P_{z_0}(\theta)d\theta.
\end{equation*}
\end{proof}
Moreover, we have part of the converse of above theorem:
\begin{theorem}\label{Poisson integral of L^1 is harmonic}
Given a function $f(\theta)\in L^1(\mathbb{T})$, its Poisson integral $P[f]$
is harmonic on $\mathbb{D}$.
\end{theorem}
\begin{proof}[Sketch of Proof]
Recall that a function $u(z)$ is said to be harmonic if 
\begin{equation*}
  \Delta u=\frac{\partial^2 u}{\partial x^2}+\frac{\partial^2 u}{\partial y^2}=0.
\end{equation*}
As Herglotz--Riesz kernel $K_z(\theta)$ is analytic over $\mathbb{D}$, 
the Poisson kernel $P_z(\theta)=\text{Re} K_z(\theta)$ is harmonic over 
$\mathbb{D}$, thus for $z\in \mathbb{D}$,
\begin{equation*}
  \Delta P[f](z)=\frac{1}{2\pi}\int \cancel{\Delta P_z(\theta)} f(\theta)d\theta=0.
\end{equation*}
\end{proof}
More generally, the same holds for Poisson integrals of measures:
\begin{theorem}
Given a finite measure $\mu$ on $\mathbb{T}$, its Poisson integral $P[\mu]$
is harmonic on $\mathbb{D}$.
\end{theorem}
\begin{proof}[Sketch of Proof]
Similar as in the previous proof, for $z\in\mathbb{D}$,
\begin{equation*}
  \Delta P[\mu](z)=\int \cancel{\Delta P_z(\theta)}\, d\mu(\theta)=0.
\end{equation*}
\end{proof}
The full converse theorem of Poisson integral formula states that the Poisson integral indeed
gives a continuous solution for Dirichlet problem, which will be later proved in Corollary 
\ref{Dirichlet problem on the disc}. Before that, we'd like to introduce more properties of 
Poisson integral.



Here are some properties for Poisson kernel:
\begin{lemma}\label{Properties of Poisson kernel}
Let $z\in\mathbb{D}$. Then
\begin{itemize}
  \item[(a)] $P_z(\theta)>0$ for all $\theta$, and
  $\displaystyle\frac{1}{2\pi}\int_0^{2\pi} P_z(\theta)\, d\theta=1$;
  \item[(b)] for any $\theta_0\in\mathbb{R}$ and $0<\delta<\pi$,
  \begin{equation*}
    \sup_{\delta\leq|\theta-\theta_0|\leq\pi}P_z(\theta)\rightarrow 0
    \quad\text{and}\quad
    \frac{1}{2\pi}\int_{\delta\leq|\theta-\theta_0|\leq\pi}P_z(\theta)\, d\theta\rightarrow 0
    \qquad (z\rightarrow e^{i\theta_0}).
  \end{equation*}
\end{itemize}
\end{lemma}

Here's an important estimate for Poisson integral of $L^p$ function:
\begin{theorem}
Let $1\leq p\leq\infty$ and let $f$ be a function in $L^p(\mathbb{T})$, and
$\|f\|_p^p:=\frac{1}{2\pi}\int_0^{2\pi}|f(\theta)|^p\, d\theta$ for $p<\infty$. Then for every
$0<r<1$,
\begin{equation*}
  \|P[f](re^{i\cdot})\|_p\leq\|f\|_p,
\end{equation*}
and hence $\sup_{0<r<1}\|P[f](re^{i\cdot})\|_p\leq\|f\|_p$.
\end{theorem}
\begin{proof}[Sketch of Proof]
If $p=\infty$, then $|P[f](z)|\leq\frac{1}{2\pi}\int P_z(\theta)|f(\theta)|\, d\theta
\leq\|f\|_\infty$.

If $1\leq p<\infty$, for $z\in\mathbb{D}$, part (a) of the lemma \ref{Properties of Poisson kernel} says that $P_z(\theta)\frac{d\theta}{2\pi}$ is a
probability measure on $\mathbb{T}$. Jensen's inequality for this probability measure gives
\begin{equation*}
  |P[f](z)|^p\leq\frac{1}{2\pi}\int_0^{2\pi} P_z(\theta)|f(\theta)|^p\, d\theta.
\end{equation*}
Set $z=re^{i\varphi}$. Note the symmetry of Poisson kernel:
\begin{equation*}
  P_{re^{i\varphi}}(\theta)=\frac{1-r^2}{1-2r\cos(\theta-\varphi)+r^2}=P_{re^{i\theta}}(\varphi),
\end{equation*}
Thus:
\begin{equation*}
  \frac{1}{2\pi}\int_0^{2\pi}|P[f](re^{i\varphi})|^p\, d\varphi
  \leq\frac{1}{2\pi}\int_0^{2\pi}|f(\theta)|^p
  \left(\frac{1}{2\pi}\int_0^{2\pi}P_{re^{i\theta}}(\varphi)\, d\varphi\right) d\theta
  =\|f\|_p^p. \qedhere
\end{equation*}
\end{proof}

The following theorem shows that the Poisson kernel acts like the identity as $r\rightarrow 1$:
\begin{theorem}\label{Poisson approximate identity}
Let $C(\mathbb{T})$ be the space of all continuous functions over $\mathbb{T}$. Then
\begin{itemize}
  \item[(a)] If $f\in C(\mathbb{T})$, then $P[f](re^{i\theta})$ converges uniformly to $f(e^{i\theta})$
  as $r\rightarrow 1$.
  \item[(b)] If $\mu$ is a finite measure on $\mathbb{T}$, then the measures
  $P[\mu](re^{i\theta})\, d\theta/(2\pi)$ converge weak-star to $\mu$ as $r\rightarrow 1$, i.e.
  \begin{equation*}
    \frac{1}{2\pi}\int_0^{2\pi} g(\theta)P[\mu](re^{i\theta})\, d\theta\rightarrow\int g\, d\mu
    \qquad (r\rightarrow 1)
  \end{equation*}
  for every $g\in C(\mathbb{T})$.
\end{itemize}
\end{theorem}
\begin{proof}[Sketch of Proof]
(a) Fix $\varphi$. By part (a) of the lemma \ref{Properties of Poisson kernel},
\begin{equation*}
  |P[f](re^{i\varphi})-f(e^{i\varphi})|
  \leq\frac{1}{2\pi}\int_{|\theta-\varphi|<\delta}P_{re^{i\varphi}}(\theta)|f(e^{i\theta})-f(e^{i\varphi})|\, d\theta
  +\frac{1}{2\pi}\int_{\delta\leq|\theta-\varphi|\leq\pi}P_{re^{i\varphi}}(\theta)|f(e^{i\theta})-f(e^{i\varphi})|\, d\theta.
\end{equation*}
As $\mathbb{T}$ is compact, $f$ is uniformly continuous on $\mathbb{T}$.
Given $\varepsilon>0$, by uniform continuity of $f$, we can pick $\delta>0$ s.t.
$|f(e^{i\theta})-f(e^{i\varphi})|<\varepsilon$ for all $\varphi$ and all $|\theta-\varphi|<\delta$,
then we have
\begin{equation*}
  \frac{1}{2\pi}\int_{|\theta-\varphi|<\delta}P_{re^{i\varphi}}(\theta)|f(e^{i\theta})-f(e^{i\varphi})|\, d\theta
  \leq \varepsilon\cdot\frac{1}{2\pi}\int P_{re^{i\varphi}}(\theta)d\theta=\varepsilon.
\end{equation*}
For this $\delta$, the
second term is at most $2\|f\|_\infty\cdot\frac{1}{2\pi}\int_{\delta\leq|\theta-\varphi|\leq\pi}
P_{re^{i\varphi}}(\theta)\, d\theta$. Since $P_{re^{i\varphi}}(\theta)=P_r(\theta-\varphi)$, this integral
does not depend on $\varphi$, and it tends to $0$ as $r\rightarrow 1$ by part (b) of the
lemma \ref{Properties of Poisson kernel}. Hence the convergence is uniform in $\varphi$.

(b) Let $g\in C(\mathbb{T})$. For fixed $r$, the Poisson kernel is bounded, so by Fubini's theorem and
the symmetry $P_{re^{i\theta}}(\varphi)=P_{re^{i\varphi}}(\theta)$,
\begin{equation*}
  \frac{1}{2\pi}\int_0^{2\pi} g(\theta)P[\mu](re^{i\theta})\, d\theta
  =\int\left(\frac{1}{2\pi}\int_0^{2\pi}P_{re^{i\varphi}}(\theta)g(\theta)\, d\theta\right) d\mu(\varphi)
  =\int P[g](re^{i\varphi})\, d\mu(\varphi).
\end{equation*}
By (a), $P[g](re^{i\varphi})\rightarrow g(e^{i\varphi})$ uniformly, and $\mu$ is finite, so the
right-hand side tends to $\int g\, d\mu$.
\end{proof}

As a consequence, the Poisson integral solves the Dirichlet problem on the disc (which is 
the converse result of the Poisson integral formula
 \ref{Poisson integral formula}):
\begin{corollary}[Dirichlet problem on the disc]\label{Dirichlet problem on the disc}
Let $f$ be a continuous function on $\mathbb{T}$. Define
\begin{equation*}
  u(z):=
  \begin{cases}
    P[f](z), & z\in\mathbb{D},\\
    f(z), & z\in\mathbb{T}.
  \end{cases}
\end{equation*}
Then $u$ is continuous on $\overline{\mathbb{D}}$ and harmonic on $\mathbb{D}$, i.e. 
a continuous solution on $\overline{\mathbb{D}}$ for the Dirichlet
problem:
\begin{equation*}
  \begin{cases}
  \Delta u=0, & z\in \mathbb{D}\\
  u=f, & z\in \mathbb{T}.
  \end{cases}
\end{equation*}
\end{corollary}
\begin{proof}
$u$ is harmonic on $\mathbb{D}$ by Theorem \ref{Poisson integral of L^1 is harmonic}. For
continuity at $e^{i\theta_0}\in\mathbb{T}$, let $z=re^{i\varphi}\in\mathbb{D}$; then
\begin{equation*}
  |u(z)-f(e^{i\theta_0})|\leq|P[f](re^{i\varphi})-f(e^{i\varphi})|+|f(e^{i\varphi})-f(e^{i\theta_0})|,
\end{equation*}
which tends to $0$ as $z\rightarrow e^{i\theta_0}$ by part (a) of
Theorem \ref{Poisson approximate identity} and the continuity of $f$.
\end{proof}

Moreover, a measure on $\mathbb{T}$ is determined by its Poisson integral:
\begin{corollary}
Let $\mu$ and $\nu$ be finite measures on $\mathbb{T}$. If $P[\mu]=P[\nu]$ on $\mathbb{D}$,
then $\mu=\nu$.
\end{corollary}
\begin{proof}
Let $\lambda:=\mu-\nu$. By linearity, $P[\lambda]=0$ on $\mathbb{D}$. For every $g\in C(\mathbb{T})$,
part (b) of Theorem \ref{Poisson approximate identity} gives
\begin{equation*}
  \int g\, d\lambda=\lim_{r\rightarrow 1}\frac{1}{2\pi}\int_0^{2\pi} g(\theta)P[\lambda](re^{i\theta})\, d\theta=0.
\end{equation*}
By the uniqueness part of the Riesz representation theorem, $\lambda=0$, thus $\mu=\nu$.
\end{proof}

Finally, Poisson integrals of $L^p$ functions and of measures are characterized by the growth
of their radial means:
\begin{theorem}\label{Characterization of Poisson integrals}
Let $u$ be a harmonic function on $\mathbb{D}$.
\begin{itemize}
  \item[(a)] If $1<p\leq\infty$, then $u=P[f]$ for some $f\in L^p(\mathbb{T})$ if and only if
  \begin{equation*}
    \sup_{0<r<1}\|u(re^{i\cdot})\|_p<\infty.
  \end{equation*}
  \item[(b)] $u=P[\mu]$ for some finite measure $\mu$ on $\mathbb{T}$ if and only if
  \begin{equation*}
    \sup_{0<r<1}\|u(re^{i\cdot})\|_1<\infty.
  \end{equation*}
  \item[(c)] (Herglotz) $u\geq 0$ on $\mathbb{D}$ if and only if $u=P[\mu]$ for some finite
  positive measure $\mu$ on $\mathbb{T}$.
\end{itemize}
\end{theorem}

\begin{proof}
(\textcolor{red}{missing})
\end{proof}

\section{Nontangential limit}
When we are discussing functions over $\mathbb{D}$, it will be good to have the proper
notion of `limit function' over the boundary $\mathbb{T}=\partial \mathbb{D}$.
This section is mainly based on \cite{rudin1974real}, 11.21 - 11.24.
\begin{definition}
Fix $\alpha>1$. For $\zeta\in\mathbb{T}$, we consider the \textbf{Stolz region}:
\begin{equation*}
  \Gamma_\alpha(\zeta):=\left\{z\in\mathbb{D}: |\zeta-z|<\alpha(1-|z|)\right\}.
\end{equation*}
A function $u:\mathbb{D}\to\mathbb{C}$ is said to have \textbf{nontangential limit} $\lambda\in\mathbb{C}$
at $\zeta\in\mathbb{T}$ if, for every $\alpha>1$,
\begin{equation*}
  \lim_{\substack{z\to\zeta\\ z\in\Gamma_\alpha(\zeta)}}u(z)=L.
\end{equation*}
When this limit exists, we write $u^*(\zeta)=\lambda$ and call it the \textbf{(nontangential) boundary value} of $u$ at $\zeta$.
This give rise to a function $u^*:\mathbb{T}\rightarrow\mathbb{C}$, for which we call the \textbf{(nontangential) boundary function}
of $u$.
\end{definition}


Before we discuss the main results related with nontangential limit, we first 
introduce a notion called \textbf{Lebesgue point}:
\begin{definition}[Lebesgue point]
Consider $(\mathbb{T},d\sigma)$, where $d\sigma := d\theta/(2\pi)$.
If $f\in L^1(\mathbb{T})$, a point $\zeta=e^{i\theta_0}\in\mathbb{T}$
is called a \textbf{Lebesgue point} of $f$ if
\begin{equation*}
  \lim_{h\downarrow 0}
  \frac{1}{2h}
  \int_{\theta_0-h}^{\theta_0+h}
  \left|f(e^{i\theta})-f(e^{i\theta_0})\right|\,d\theta=0.
\end{equation*}
\end{definition}

\begin{remark}
If $f$ is continuous at $\xi$, then $\xi$ is a Lebesgue point of $f$. So one can
think of Lebesgue points of a function are those points where in average sense, the
function does not oscillate too much.
\end{remark}

\begin{theorem}[Lebesgue differentiation theorem]
\label{thm:lebesgue_differentiation}
Let $f\in L^1(\mathbb{T})$. Then almost every point of $\mathbb{T}$
 is a Lebesgue point of $f$.
\end{theorem}

The following result tells us that nontangential limit of
the Poisson integral of a function gives back the function value itself
at all lebesgue points of that function:
\begin{theorem}
\label{thm:poisson_nontangential_lebesgue_point}
Let $f\in L^1(\mathbb{T})$ be real-valued. If $\zeta\in\mathbb{T}$ is a
Lebesgue point of $f$, then the Poisson integral $P[f]$ has nontangential
limit $f(\zeta)$ at $\zeta$. That is, for every $\alpha>1$,
\begin{equation*}
  \lim_{\substack{z\to\zeta\\ z\in\Gamma_\alpha(\zeta)}}P[f](z)=f(\zeta).
\end{equation*}
\end{theorem}

Then using Lebesgue differentiation theorem \ref{thm:lebesgue_differentiation},
we get an a.e. recovery of boundary data:
\begin{corollary}[Almost everywhere recovery of boundary data]
\label{cor:poisson_nontangential_ae}
Let $f\in L^1(\mathbb{T})$ be real-valued. Then $P[f]$ has nontangential
limits almost everywhere on $\mathbb{T}$, and
\begin{equation*}
  (P[f])^*(\zeta)=f(\zeta)
  \qquad\text{for almost every }\zeta\in\mathbb{T}.
\end{equation*}
\end{corollary}

\begin{theorem}
\label{thm:singular_poisson_nontangential_zero}
Let $\nu$ be a finite measure on $\mathbb{T}$ s.t. $\nu\bot d\sigma$, i.e. singular
with respect to $d\sigma = d\theta/(2\pi)$. Then its Poisson integral $P[\nu]$ has
nontangential limit 0 at almost every point of $\mathbb{T}$. For every $\alpha>1$,
\begin{equation*}
  \lim_{\substack{z\to\zeta\\ z\in\Gamma_\alpha(\zeta)}}P[\nu](z)=0
  \qquad\text{for a.e. }\zeta\in\mathbb{T}.
\end{equation*}
\end{theorem}

\begin{theorem}[Fatou's theorem for finite measures]
\label{thm:fatou_signed_measure}
Let $\mu$ be a finite measure on $\mathbb{T}$ with Lebesgue
decomposition
\begin{equation*}
  d\mu=f\,\frac{d\theta}{2\pi}+d\mu_s,
  \qquad f\in L^1(\mathbb{T}),
  \qquad \mu_s\perp\frac{d\theta}{2\pi}.
\end{equation*}
Then $P[\mu]$ has nontangential limits a.e. on $\mathbb{T}$, and for every $\alpha>1$,
\begin{equation*}
  \lim_{\substack{z\to\zeta\\ z\in\Gamma_\alpha(\zeta)}}P[\mu](z)=f(\zeta)
  \qquad\text{for a.e. }\zeta\in\mathbb{T}.
\end{equation*}
\end{theorem}
\begin{proof}
By linearity, $P[\mu]=P[f]+P[\mu_s]$. Apply
Corollary~\ref{cor:poisson_nontangential_ae} and
Theorem~\ref{thm:singular_poisson_nontangential_zero} will get the desired result.
\end{proof}





\section{Subharmonic function}
\label{subsec:subharmonic}
To investigate the main objects in this thesis,
which are Hardy spaces and Nevanlinna class, it will be good to know the properties of 
function like $|f|^p$, $\log |f|$ and $\log^+ |f|$
for analytic function $f$. These functions belong to an important class of functions called 
\textbf{subharmonic functions}. This section is based on \cite{garnett2006bounded}, I.6.
\begin{definition}
Let $\Omega$ be an open set in the plane. A \textbf{subharmonic function} on $\Omega$ is a function
$v : \Omega \to [-\infty, \infty)$ such that
\begin{itemize}
  \item[(a)] $v$ is \emph{upper semicontinous}, i.e. $v(z_0)\geq \lim_{z\rightarrow z_0}v(z)$ 
  for all $z_0\in\Omega$;
  \item[(b)] for each $z_0 \in \Omega$ there is $r(z_0) > 0$ such that the disc 
$B(z_0, r(z_0)) = \{z : |z - z_0| < r(z_0)\}$ is contained in $\Omega$ and 
such that for every $r < r(z_0)$,
\begin{equation}\label{MVP of subharmonic functions}
  v(z_0) \leq \frac{1}{\pi r^2} \iint_{B(z_0,r)} v(z)\, dx\, dy.
\end{equation}
We often call (\ref{MVP of subharmonic functions}) the `\textbf{submean value property}'.
\end{itemize}
\end{definition}
\begin{remark}
Note that the integral in \ref{MVP of subharmonic functions} either
converges or diverges to $-\infty$, due to the upper-semicontinuity by the following
lemma:
\begin{lemma}\label{boundedness of uppersemicont fct}
Let $X$ be a topological space.
If $v:X\rightarrow [-\infty,\infty)$ is upper-semicontinuous, then $v$ is bounded
above on each compact subset $K\subset X$.
\end{lemma}
\end{remark}

Here are some examples of subharmonic functions:
\begin{example}
Every harmonic function is subharmonic, due to the \textbf{mean value property} of harmonic function:
\begin{theorem}
Let $\Omega$ be an open set in the plane and let $u$ be a harmonic function on $\Omega$. Then for
every $z_0\in\Omega$ and every $r>0$ such that the closed disc $\overline{B(z_0,r)}$ is contained
in $\Omega$,
\begin{equation*}
  u(z_0) = \frac{1}{2\pi}\int_0^{2\pi} u(z_0+re^{i\theta})\, d\theta
  = \frac{1}{\pi r^2} \iint_{B(z_0,r)} u(z)\, dx\, dy.
\end{equation*}
\end{theorem}
\end{example}
\begin{example}\label{log |f| is subharmonic}
Let $f(z)$ be an analytic function on a domain $\Omega\subset \mathbb{C}$, then 
$v(z)=\log |f(z)|$ is subharmonic on $\Omega$: 
\begin{itemize}
  \item $v$ is upper semicontinuous;
  \item For $z_0$ s.t. $f(z_0)=0$, condition \ref{MVP of subharmonic functions} trivially
  satisfied;
  \item For $z_0$ s.t. $f(z_0)\neq 0$, we can write $v(z)=\text{Re}(\log f(z))$ on some neighborhood
  of $z_0$, which is harmonic on this neighborhood.
\end{itemize}
\end{example}
\begin{example}
Let $f(z)$ be an analytic function on a domain $\Omega\subset \mathbb{C}$, then 
$|f(z)|^p$ is subharmonic on $\Omega$ as 
\begin{equation*}
  |f(z)|^p=\exp(p \log |f(z)|)
\end{equation*}
So the condition \ref{MVP of subharmonic functions} can be proved by the fact that
$\log |f(z)|$ is subharmonic and Jensen's inequality as $\exp$ is convex. More generally, we
have the following lemma:
\begin{lemma}
Given a subharmonic funciton $v$ on $\Omega$ and an increasing convex function $\varphi$ on
$[-\infty, \infty)$ which is continuous at $-\infty$, then $\varphi\circ v$ is subharmonic on
$\Omega$.
\end{lemma}
With this lemma, we can see that
\begin{equation*}
    \log^+ |f(z)|=\max(\log |f(z)|,0)
\end{equation*}
is also a subharmonic function.
\end{example}

There is an equivalent characterization of subharmonicity:
\begin{theorem}[Harmonic comparison characterization]\label{Harmonic comparison characterization}
Let $u:\Omega\rightarrow [-\infty,\infty)$ be an upper semicontinuous function. Then $u$ is 
subharmonic on $\Omega$ if and only if the following condition holds: If $v$ is a continuous function
on $\overline{B(z_0,r)}\subset \Omega$ that is harmonic on $B(z_0,r)$ s.t. for all 
$\xi\in \partial B(z_0,r)$, we have $u(\xi)\leq v(\xi)$, then 
\begin{equation*}
  u(z)\leq v(z),\quad z\in B(z_0,r).
\end{equation*}
\end{theorem}
\begin{proof}
(\textcolor{red}{missing})
\end{proof}

\begin{corollary}
If $\Omega$ is a connected open set and $v$ is a subharmonic function on $\Omega$ s.t. $v\not\equiv -\infty$,
then given any closed disc $\overline{B(z_0,r)}\subset \Omega$, 
\begin{equation*}
    \frac{1}{2\pi}\int v(z_0+r e^{i\theta})d\theta \neq -\infty.
\end{equation*}
\end{corollary}
\begin{proof}[Sketch of Proof]
First recall a lemma:
\begin{lemma}
Let $X$ be a metric space, $v$ be an upper-semicontinuous function. Then we can always find
a sequence $u_n$ of continuous function that decreasing to $v$ on $X$.
\end{lemma}
Now we pick $u_n$ to be a sequence of continuous functions decreasing to $v$ on 
$\partial B(z_0,r)$. Let $U_n(z)=P[u_n](z)$, i.e. the harmonic extension of $u_n$ to $B(z_0,r)$.
Assume for the sake of contradiction that we can find some $\overline{B(z_0,r)}\subset\Omega$ s.t.
\begin{equation*}
  \frac{1}{2\pi}\int v(z_0+r e^{i\theta})d\theta = -\infty.
\end{equation*}
Recall Lemma \ref{boundedness of uppersemicont fct}, $v$ is bounded above by some
constant $M\geq 0$. Write $v=v^+-v^-$ with $v^+=\max(v,0)$ and $v^-=\max(-v,0)$, so that
$0\leq v^+\leq M$ and $0\leq v^-\leq\infty$ on $\partial B(z_0,r)$. Since
\begin{equation*}
  \frac{1}{2\pi}\int v^+(z_0+re^{i\theta})\, d\theta\leq M<\infty,
\end{equation*}
the assumption is equivalent to
\begin{equation*}
  \frac{1}{2\pi}\int v^-(z_0+re^{i\theta})\, d\theta=\infty.
\end{equation*}
Now fix $z$ with $|z|<1$. Since $1-|z|\leq |e^{i\theta}-z|\leq 1+|z|$, we have
\begin{equation*}
  0<\frac{1-|z|}{1+|z|}\leq P_z(\theta)=\frac{1-|z|^2}{|e^{i\theta}-z|^2}\leq\frac{1+|z|}{1-|z|}
  \quad\text{for all }\theta.
\end{equation*}
Therefore
\begin{align*}
  \frac{1}{2\pi}\int P_z(\theta)\, v^+(z_0+re^{i\theta})\, d\theta
  &\leq \frac{1+|z|}{1-|z|}\, M<\infty,\\
  \frac{1}{2\pi}\int P_z(\theta)\, v^-(z_0+re^{i\theta})\, d\theta
  &\geq \frac{1-|z|}{1+|z|}\cdot\frac{1}{2\pi}\int v^-(z_0+re^{i\theta})\, d\theta=\infty.
\end{align*}
Therefore,
\begin{equation*}
  \frac{1}{2\pi}\int P_z(\theta)\, v(z_0+re^{i\theta})\, d\theta = -\infty,\quad |z|<1.
\end{equation*}
As
\begin{equation*}
  U_n(z)=\frac{1}{2\pi}\int P_z(\theta)\, u_n(z_0+re^{i\theta})\, d\theta 
\end{equation*}
and $u_n\downarrow v$, so $U_n(z)\rightarrow -\infty$ for each $z\in B(z_0,r)$. By Harmonic
comparison \ref{Harmonic comparison characterization}, we know that $U_n \geq v$ over $B(z_0,r)$,
so $v=-\infty$ on $B(z_0,r)$. With above result, we can actually extend this local property globally
using standard topological argument as follow: consider 
\begin{equation*}
  A:=\{z\in \Omega: v(z)\equiv -\infty\text{ on a neighborhood of }z\}
\end{equation*}
then one can show that it is open and closed. As $z_0\in A$, it is nonempty, thus $v\equiv -\infty$
over $\Omega$, which gives a contradiction.

\end{proof}

To investigate further properties of subharmonic functions, we'd like to first define
an auxiliary function called \textbf{Poisson modification}:
\begin{definition}[Poisson modification]\label{Poisson modification}
Let $v:\mathbb{D}\rightarrow[-\infty,\infty)$ be an upper semicontinuous function and let
$0<r<1$. The \textbf{Poisson modification} of $v$ on $B(0,r)$ is the function
\begin{equation*}
  v_r(z):=
  \begin{cases}
    v(z), & r\leq|z|<1,\\[4pt]
    \displaystyle\frac{1}{2\pi}\int P_{z/r}(\theta)\, v(re^{i\theta})\, d\theta, & |z|<r.
  \end{cases}
\end{equation*}
\end{definition}
By Lemma \ref{boundedness of uppersemicont fct}, $v$ is bounded above on $\partial B(0,r)$, so
the integral is well-defined with values in $[-\infty,\infty)$.

\begin{theorem}\label{Properties of Poisson modification}
Let $v$ be a subharmonic function on $\mathbb{D}$ s.t. $v\not\equiv-\infty$, and let $0<r<1$.
Then
\begin{itemize}
  \item[(a)] $v_r$ is finite and harmonic on $B(0,r)$;
  \item[(b)] $v_r$ is subharmonic on $\mathbb{D}$;
  \item[(c)] $v(z)\leq v_r(z)$ for all $z\in\mathbb{D}$;
  \item[(d)] $v_r(z)$ is an increasing function of $r$.
\end{itemize}
\end{theorem}
\begin{proof}
(\textcolor{red}{missing})
\end{proof}
If we take $z=0$ in Theorem \ref{Properties of Poisson modification}(d), we get:
\begin{corollary}
If $v(z)$ is a subharmonic function on $\mathbb{D}$, then
\begin{equation*}
  m(r)=v_r(0)=\frac{1}{2\pi}\int v(r e^{i\theta})d\theta
\end{equation*}
is an increasing function of $r\in(0,1)$.
\end{corollary}

\begin{definition}[Harmonic majorant]\label{Harmonic majorant}
Let $v$ be a subharmonic function on an open set $\Omega$. A \textbf{harmonic majorant} of $v$
is a harmonic function $U$ on $\Omega$ s.t. $v(z)\leq U(z)$ for all $z\in\Omega$. A harmonic
majorant $u$ of $v$ is called the \textbf{least harmonic majorant} of $v$ if $u(z)\leq U(z)$,
$z\in\Omega$, for every harmonic majorant $U$ of $v$.
\end{definition}

\begin{theorem}\label{Harmonic majorant criterion}
Let $v$ be a subharmonic function on $\mathbb{D}$ s.t. $v\not\equiv-\infty$. Then $v$ has a
harmonic majorant if and only if
\begin{equation*}
  \sup_{0<r<1}\frac{1}{2\pi}\int v(re^{i\theta})\, d\theta=\sup_{0<r<1}v_r(0)<\infty.
\end{equation*}
In this case, the least harmonic majorant of $v$ is
\begin{equation*}
  u(z)=\lim_{r\rightarrow 1}\frac{1}{2\pi}\int P_{z/r}(\theta)\, v(re^{i\theta})\, d\theta
  =\lim_{r\rightarrow 1}v_r(z),\quad z\in\mathbb{D}.
\end{equation*}
\end{theorem}

\begin{proof}
\textcolor{red}{missing}
\end{proof}.

\chapter{Background}
\label{sec:background}
Before we talk about the main object in this thesis, we need to first introduce some analysis
backgrounds.
\section{Poisson integral}
Let $\mathbb{D}:=\{z\in\C: |z|<1\}$ and $\mathbb{T}=\partial \mathbb{D}$. 
When given a continuous function $f$ on $\mathbb{T}$, the famous \textbf{Dirichlet
problem}, which is to find the solution for the following equations:
\begin{equation*}
  \begin{cases}
  \Delta u=0, & z\in \mathbb{D}\\
  u=f, & z\in \mathbb{T}.
  \end{cases}
\end{equation*}
is asking for a harmonic extension $u$ over $\mathbb{D}$ with the boundary data $f$ on $\mathbb{T}$.
To give a solution for this problem, we'd like to introduce the notion called \textbf{Poisson integral}.
This section is based on \cite{garnett2006bounded}, I.3.
\begin{definition}
Let $c\in\mathbb{C}$. The \textbf{Herglotz--Riesz kernel centered at $c$} is the function
\begin{equation*}
  K(c;z,w) := \frac{(w-c)+(z-c)}{(w-c)-(z-c)} = \frac{w+z-2c}{w-z},
  \qquad z,w\in\mathbb{C},\ w\neq z.
\end{equation*}
The \textbf{Poisson kernel centered at $c$} is defined as
\begin{equation*}
  P(c;z,w) := \text{Re}\; K(c;z,w)= \frac{|w-c|^2-|z-c|^2}{|w-z|^2},
  \qquad z,w\in\mathbb{C},\ w\neq z.
\end{equation*}
and we denote $K(z,w):= K(0;z,w)$, $P(z,w):= P(0;z,w)$, and $K_z(\theta):= K(z,e^{i\theta})$,
 $P_z(\theta):= P(z,e^{i\theta})$. 
\end{definition}

\begin{definition}
Given a function $f(\theta)\in L^1(\mathbb{T})$, we define the \textbf{Poisson integral} of $f$ 
to be the function
\begin{equation*}
  P[f](z):=\frac{1}{2\pi}\int P_z(\theta) f(\theta)d\theta.
\end{equation*}
In general, given a finite (signed) measure $\mu$ on $\mathbb{T}$, we define the \textbf{Poisson integral} of $\mu$
to be the function
\begin{equation*}
  P[\mu](z):=\int P_z(\theta) d\mu (\theta).
\end{equation*}
The first definition is a special case of the second one if we set 
$d\mu(\theta)=\frac{1}{2\pi} f(\theta)d\theta$.
\end{definition}

\begin{theorem}[Poisson integral formula]\label{Poisson integral formula}
Let $u$ be a continuous function on the closed unit disc $\overline{\mathbb{D}}$. If $u$ is
harmonic on $\mathbb{D}$, then
\begin{equation*}
  u(z)=\frac{1}{2\pi}\int_0^{2\pi} P_z(\theta)\, u(e^{i\theta})\, d\theta
  \qquad \text{for all } z\in\mathbb{D}.
\end{equation*}
\end{theorem}
\begin{proof}
Recall that if $u$ is harmonic on $\mathbb{D}$, then it has the 
\textbf{mean value property}: 
\begin{equation*}
  u(0)=\frac{1}{2\pi}\int_0^{2\pi} u(e^{i\theta})d\theta
\end{equation*}
Now given $z_0=re^{i\theta_0}$, consider the M\"{o}bius transformation 
$\tau(z)=(z-z_0)/(1-\bar{z}_0 z)$, and write $\tau(e^{i\theta})=e^{i\varphi}$.
Then we can compute that $d\varphi/d\theta=P_{z_0}(\theta)$, therefore
\begin{equation*}
  u(z_0)=u(\tau^{-1}(0))=\frac{1}{2\pi}\int_0^{2\pi} (u\circ\tau^{-1})(e^{i\varphi})
  d\varphi=\frac{1}{2\pi}\int_0^{2\pi} u(e^{i\theta})P_{z_0}(\theta)d\theta.
\end{equation*}
\end{proof}
Moreover, we have part of the converse of above theorem:
\begin{theorem}\label{Poisson integral of L^1 is harmonic}
Given a function $f(\theta)\in L^1(\mathbb{T})$, its Poisson integral $P[f]$
is harmonic on $\mathbb{D}$.
\end{theorem}
\begin{proof}[Sketch of Proof]
Recall that a function $u(z)$ is said to be harmonic if 
\begin{equation*}
  \Delta u=\frac{\partial^2 u}{\partial x^2}+\frac{\partial^2 u}{\partial y^2}=0.
\end{equation*}
As Herglotz--Riesz kernel $K_z(\theta)$ is analytic over $\mathbb{D}$, 
the Poisson kernel $P_z(\theta)=\text{Re} K_z(\theta)$ is harmonic over 
$\mathbb{D}$, thus for $z\in \mathbb{D}$,
\begin{equation*}
  \Delta P[f](z)=\frac{1}{2\pi}\int \cancel{\Delta P_z(\theta)} f(\theta)d\theta=0.
\end{equation*}
\end{proof}
More generally, the same holds for Poisson integrals of measures:
\begin{theorem}
Given a finite measure $\mu$ on $\mathbb{T}$, its Poisson integral $P[\mu]$
is harmonic on $\mathbb{D}$.
\end{theorem}
\begin{proof}[Sketch of Proof]
Similar as in the previous proof, for $z\in\mathbb{D}$,
\begin{equation*}
  \Delta P[\mu](z)=\int \cancel{\Delta P_z(\theta)}\, d\mu(\theta)=0.
\end{equation*}
\end{proof}
The full converse theorem of Poisson integral formula states that the Poisson integral indeed
gives a continuous solution for Dirichlet problem, which will be later proved in Corollary 
\ref{Dirichlet problem on the disc}. Before that, we'd like to introduce more properties of 
Poisson integral.



Here are some properties for Poisson kernel:
\begin{lemma}\label{Properties of Poisson kernel}
Let $z\in\mathbb{D}$. Then
\begin{itemize}
  \item[(a)] $P_z(\theta)>0$ for all $\theta$, and
  $\displaystyle\frac{1}{2\pi}\int_0^{2\pi} P_z(\theta)\, d\theta=1$;
  \item[(b)] for any $\theta_0\in\mathbb{R}$ and $0<\delta<\pi$,
  \begin{equation*}
    \sup_{\delta\leq|\theta-\theta_0|\leq\pi}P_z(\theta)\rightarrow 0
    \quad\text{and}\quad
    \frac{1}{2\pi}\int_{\delta\leq|\theta-\theta_0|\leq\pi}P_z(\theta)\, d\theta\rightarrow 0
    \qquad (z\rightarrow e^{i\theta_0}).
  \end{equation*}
\end{itemize}
\end{lemma}

Here's an important estimate for Poisson integral of $L^p$ function:
\begin{theorem}
Let $1\leq p\leq\infty$ and let $f$ be a function in $L^p(\mathbb{T})$, and
$\|f\|_p^p:=\frac{1}{2\pi}\int_0^{2\pi}|f(\theta)|^p\, d\theta$ for $p<\infty$. Then for every
$0<r<1$,
\begin{equation*}
  \|P[f](re^{i\cdot})\|_p\leq\|f\|_p,
\end{equation*}
and hence $\sup_{0<r<1}\|P[f](re^{i\cdot})\|_p\leq\|f\|_p$.
\end{theorem}
\begin{proof}[Sketch of Proof]
If $p=\infty$, then $|P[f](z)|\leq\frac{1}{2\pi}\int P_z(\theta)|f(\theta)|\, d\theta
\leq\|f\|_\infty$.

If $1\leq p<\infty$, for $z\in\mathbb{D}$, part (a) of the lemma \ref{Properties of Poisson kernel} says that $P_z(\theta)\frac{d\theta}{2\pi}$ is a
probability measure on $\mathbb{T}$. Jensen's inequality for this probability measure gives
\begin{equation*}
  |P[f](z)|^p\leq\frac{1}{2\pi}\int_0^{2\pi} P_z(\theta)|f(\theta)|^p\, d\theta.
\end{equation*}
Set $z=re^{i\varphi}$. Note the symmetry of Poisson kernel:
\begin{equation*}
  P_{re^{i\varphi}}(\theta)=\frac{1-r^2}{1-2r\cos(\theta-\varphi)+r^2}=P_{re^{i\theta}}(\varphi),
\end{equation*}
Thus:
\begin{equation*}
  \frac{1}{2\pi}\int_0^{2\pi}|P[f](re^{i\varphi})|^p\, d\varphi
  \leq\frac{1}{2\pi}\int_0^{2\pi}|f(\theta)|^p
  \left(\frac{1}{2\pi}\int_0^{2\pi}P_{re^{i\theta}}(\varphi)\, d\varphi\right) d\theta
  =\|f\|_p^p. \qedhere
\end{equation*}
\end{proof}

The following theorem shows that the Poisson kernel acts like the identity as $r\rightarrow 1$:
\begin{theorem}\label{Poisson approximate identity}
Let $C(\mathbb{T})$ be the space of all continuous functions over $\mathbb{T}$. Then
\begin{itemize}
  \item[(a)] If $f\in C(\mathbb{T})$, then $P[f](re^{i\theta})$ converges uniformly to $f(e^{i\theta})$
  as $r\rightarrow 1$.
  \item[(b)] If $\mu$ is a finite measure on $\mathbb{T}$, then the measures
  $P[\mu](re^{i\theta})\, d\theta/(2\pi)$ converge weak-star to $\mu$ as $r\rightarrow 1$, i.e.
  \begin{equation*}
    \frac{1}{2\pi}\int_0^{2\pi} g(\theta)P[\mu](re^{i\theta})\, d\theta\rightarrow\int g\, d\mu
    \qquad (r\rightarrow 1)
  \end{equation*}
  for every $g\in C(\mathbb{T})$.
\end{itemize}
\end{theorem}
\begin{proof}[Sketch of Proof]
(a) Fix $\varphi$. By part (a) of the lemma \ref{Properties of Poisson kernel},
\begin{equation*}
  |P[f](re^{i\varphi})-f(e^{i\varphi})|
  \leq\frac{1}{2\pi}\int_{|\theta-\varphi|<\delta}P_{re^{i\varphi}}(\theta)|f(e^{i\theta})-f(e^{i\varphi})|\, d\theta
  +\frac{1}{2\pi}\int_{\delta\leq|\theta-\varphi|\leq\pi}P_{re^{i\varphi}}(\theta)|f(e^{i\theta})-f(e^{i\varphi})|\, d\theta.
\end{equation*}
As $\mathbb{T}$ is compact, $f$ is uniformly continuous on $\mathbb{T}$.
Given $\varepsilon>0$, by uniform continuity of $f$, we can pick $\delta>0$ s.t.
$|f(e^{i\theta})-f(e^{i\varphi})|<\varepsilon$ for all $\varphi$ and all $|\theta-\varphi|<\delta$,
then we have
\begin{equation*}
  \frac{1}{2\pi}\int_{|\theta-\varphi|<\delta}P_{re^{i\varphi}}(\theta)|f(e^{i\theta})-f(e^{i\varphi})|\, d\theta
  \leq \varepsilon\cdot\frac{1}{2\pi}\int P_{re^{i\varphi}}(\theta)d\theta=\varepsilon.
\end{equation*}
For this $\delta$, the
second term is at most $2\|f\|_\infty\cdot\frac{1}{2\pi}\int_{\delta\leq|\theta-\varphi|\leq\pi}
P_{re^{i\varphi}}(\theta)\, d\theta$. Since $P_{re^{i\varphi}}(\theta)=P_r(\theta-\varphi)$, this integral
does not depend on $\varphi$, and it tends to $0$ as $r\rightarrow 1$ by part (b) of the
lemma \ref{Properties of Poisson kernel}. Hence the convergence is uniform in $\varphi$.

(b) Let $g\in C(\mathbb{T})$. For fixed $r$, the Poisson kernel is bounded, so by Fubini's theorem and
the symmetry $P_{re^{i\theta}}(\varphi)=P_{re^{i\varphi}}(\theta)$,
\begin{equation*}
  \frac{1}{2\pi}\int_0^{2\pi} g(\theta)P[\mu](re^{i\theta})\, d\theta
  =\int\left(\frac{1}{2\pi}\int_0^{2\pi}P_{re^{i\varphi}}(\theta)g(\theta)\, d\theta\right) d\mu(\varphi)
  =\int P[g](re^{i\varphi})\, d\mu(\varphi).
\end{equation*}
By (a), $P[g](re^{i\varphi})\rightarrow g(e^{i\varphi})$ uniformly, and $\mu$ is finite, so the
right-hand side tends to $\int g\, d\mu$.
\end{proof}

As a consequence, the Poisson integral solves the Dirichlet problem on the disc (which is 
the converse result of the Poisson integral formula
 \ref{Poisson integral formula}):
\begin{corollary}[Dirichlet problem on the disc]\label{Dirichlet problem on the disc}
Let $f$ be a continuous function on $\mathbb{T}$. Define
\begin{equation*}
  u(z):=
  \begin{cases}
    P[f](z), & z\in\mathbb{D},\\
    f(z), & z\in\mathbb{T}.
  \end{cases}
\end{equation*}
Then $u$ is continuous on $\overline{\mathbb{D}}$ and harmonic on $\mathbb{D}$, i.e. 
a continuous solution on $\overline{\mathbb{D}}$ for the Dirichlet
problem:
\begin{equation*}
  \begin{cases}
  \Delta u=0, & z\in \mathbb{D}\\
  u=f, & z\in \mathbb{T}.
  \end{cases}
\end{equation*}
\end{corollary}
\begin{proof}
$u$ is harmonic on $\mathbb{D}$ by Theorem \ref{Poisson integral of L^1 is harmonic}. For
continuity at $e^{i\theta_0}\in\mathbb{T}$, let $z=re^{i\varphi}\in\mathbb{D}$; then
\begin{equation*}
  |u(z)-f(e^{i\theta_0})|\leq|P[f](re^{i\varphi})-f(e^{i\varphi})|+|f(e^{i\varphi})-f(e^{i\theta_0})|,
\end{equation*}
which tends to $0$ as $z\rightarrow e^{i\theta_0}$ by part (a) of
Theorem \ref{Poisson approximate identity} and the continuity of $f$.
\end{proof}

Moreover, a measure on $\mathbb{T}$ is determined by its Poisson integral:
\begin{corollary}
Let $\mu$ and $\nu$ be finite measures on $\mathbb{T}$. If $P[\mu]=P[\nu]$ on $\mathbb{D}$,
then $\mu=\nu$.
\end{corollary}
\begin{proof}
Let $\lambda:=\mu-\nu$. By linearity, $P[\lambda]=0$ on $\mathbb{D}$. For every $g\in C(\mathbb{T})$,
part (b) of Theorem \ref{Poisson approximate identity} gives
\begin{equation*}
  \int g\, d\lambda=\lim_{r\rightarrow 1}\frac{1}{2\pi}\int_0^{2\pi} g(\theta)P[\lambda](re^{i\theta})\, d\theta=0.
\end{equation*}
By the uniqueness part of the Riesz representation theorem, $\lambda=0$, thus $\mu=\nu$.
\end{proof}

Finally, Poisson integrals of $L^p$ functions and of measures are characterized by the growth
of their radial means:
\begin{theorem}\label{Characterization of Poisson integrals}
Let $u$ be a harmonic function on $\mathbb{D}$.
\begin{itemize}
  \item[(a)] If $1<p\leq\infty$, then $u=P[f]$ for some $f\in L^p(\mathbb{T})$ if and only if
  \begin{equation*}
    \sup_{0<r<1}\|u(re^{i\cdot})\|_p<\infty.
  \end{equation*}
  \item[(b)] $u=P[\mu]$ for some finite measure $\mu$ on $\mathbb{T}$ if and only if
  \begin{equation*}
    \sup_{0<r<1}\|u(re^{i\cdot})\|_1<\infty.
  \end{equation*}
  \item[(c)] (Herglotz) $u\geq 0$ on $\mathbb{D}$ if and only if $u=P[\mu]$ for some finite
  positive measure $\mu$ on $\mathbb{T}$.
\end{itemize}
\end{theorem}

\begin{proof}
(\textcolor{red}{missing})
\end{proof}

\section{Nontangential limit}
When we are discussing functions over $\mathbb{D}$, it will be good to have the proper
notion of `limit function' over the boundary $\mathbb{T}=\partial \mathbb{D}$.
This section is mainly based on \cite{rudin1974real}, 11.21 - 11.24.
\begin{definition}
Fix $\alpha>1$. For $\zeta\in\mathbb{T}$, we consider the \textbf{Stolz region}:
\begin{equation*}
  \Gamma_\alpha(\zeta):=\left\{z\in\mathbb{D}: |\zeta-z|<\alpha(1-|z|)\right\}.
\end{equation*}
A function $u:\mathbb{D}\to\mathbb{C}$ is said to have \textbf{nontangential limit} $\lambda\in\mathbb{C}$
at $\zeta\in\mathbb{T}$ if, for every $\alpha>1$,
\begin{equation*}
  \lim_{\substack{z\to\zeta\\ z\in\Gamma_\alpha(\zeta)}}u(z)=L.
\end{equation*}
When this limit exists, we write $u^*(\zeta)=\lambda$ and call it the \textbf{(nontangential) boundary value} of $u$ at $\zeta$.
This give rise to a function $u^*:\mathbb{T}\rightarrow\mathbb{C}$, for which we call the \textbf{(nontangential) boundary function}
of $u$.
\end{definition}


Before we discuss the main results related with nontangential limit, we first 
introduce a notion called \textbf{Lebesgue point}:
\begin{definition}[Lebesgue point]
Consider $(\mathbb{T},d\sigma)$, where $d\sigma := d\theta/(2\pi)$.
If $f\in L^1(\mathbb{T})$, a point $\zeta=e^{i\theta_0}\in\mathbb{T}$
is called a \textbf{Lebesgue point} of $f$ if
\begin{equation*}
  \lim_{h\downarrow 0}
  \frac{1}{2h}
  \int_{\theta_0-h}^{\theta_0+h}
  \left|f(e^{i\theta})-f(e^{i\theta_0})\right|\,d\theta=0.
\end{equation*}
\end{definition}

\begin{remark}
If $f$ is continuous at $\xi$, then $\xi$ is a Lebesgue point of $f$. So one can
think of Lebesgue points of a function are those points where in average sense, the
function does not oscillate too much.
\end{remark}

\begin{theorem}[Lebesgue differentiation theorem]
\label{thm:lebesgue_differentiation}
Let $f\in L^1(\mathbb{T})$. Then almost every point of $\mathbb{T}$
 is a Lebesgue point of $f$.
\end{theorem}

The following result tells us that nontangential limit of
the Poisson integral of a function gives back the function value itself
at all lebesgue points of that function:
\begin{theorem}
\label{thm:poisson_nontangential_lebesgue_point}
Let $f\in L^1(\mathbb{T})$ be real-valued. If $\zeta\in\mathbb{T}$ is a
Lebesgue point of $f$, then the Poisson integral $P[f]$ has nontangential
limit $f(\zeta)$ at $\zeta$. That is, for every $\alpha>1$,
\begin{equation*}
  \lim_{\substack{z\to\zeta\\ z\in\Gamma_\alpha(\zeta)}}P[f](z)=f(\zeta).
\end{equation*}
\end{theorem}

Then using Lebesgue differentiation theorem \ref{thm:lebesgue_differentiation},
we get an a.e. recovery of boundary data:
\begin{corollary}[Almost everywhere recovery of boundary data]
\label{cor:poisson_nontangential_ae}
Let $f\in L^1(\mathbb{T})$ be real-valued. Then $P[f]$ has nontangential
limits almost everywhere on $\mathbb{T}$, and
\begin{equation*}
  (P[f])^*(\zeta)=f(\zeta)
  \qquad\text{for almost every }\zeta\in\mathbb{T}.
\end{equation*}
\end{corollary}

\begin{theorem}
\label{thm:singular_poisson_nontangential_zero}
Let $\nu$ be a finite measure on $\mathbb{T}$ s.t. $\nu\bot d\sigma$, i.e. singular
with respect to $d\sigma = d\theta/(2\pi)$. Then its Poisson integral $P[\nu]$ has
nontangential limit 0 at almost every point of $\mathbb{T}$. For every $\alpha>1$,
\begin{equation*}
  \lim_{\substack{z\to\zeta\\ z\in\Gamma_\alpha(\zeta)}}P[\nu](z)=0
  \qquad\text{for a.e. }\zeta\in\mathbb{T}.
\end{equation*}
\end{theorem}

\begin{theorem}[Fatou's theorem for finite measures]
\label{thm:fatou_signed_measure}
Let $\mu$ be a finite measure on $\mathbb{T}$ with Lebesgue
decomposition
\begin{equation*}
  d\mu=f\,\frac{d\theta}{2\pi}+d\mu_s,
  \qquad f\in L^1(\mathbb{T}),
  \qquad \mu_s\perp\frac{d\theta}{2\pi}.
\end{equation*}
Then $P[\mu]$ has nontangential limits a.e. on $\mathbb{T}$, and for every $\alpha>1$,
\begin{equation*}
  \lim_{\substack{z\to\zeta\\ z\in\Gamma_\alpha(\zeta)}}P[\mu](z)=f(\zeta)
  \qquad\text{for a.e. }\zeta\in\mathbb{T}.
\end{equation*}
\end{theorem}
\begin{proof}
By linearity, $P[\mu]=P[f]+P[\mu_s]$. Apply
Corollary~\ref{cor:poisson_nontangential_ae} and
Theorem~\ref{thm:singular_poisson_nontangential_zero} will get the desired result.
\end{proof}





\section{Subharmonic function}
\label{subsec:subharmonic}
To investigate the main objects in this thesis,
which are Hardy spaces and Nevanlinna class, it will be good to know the properties of 
function like $|f|^p$, $\log |f|$ and $\log^+ |f|$
for analytic function $f$. These functions belong to an important class of functions called 
\textbf{subharmonic functions}. This section is based on \cite{garnett2006bounded}, I.6.
\begin{definition}
Let $\Omega$ be an open set in the plane. A \textbf{subharmonic function} on $\Omega$ is a function
$v : \Omega \to [-\infty, \infty)$ such that
\begin{itemize}
  \item[(a)] $v$ is \emph{upper semicontinous}, i.e. $v(z_0)\geq \lim_{z\rightarrow z_0}v(z)$ 
  for all $z_0\in\Omega$;
  \item[(b)] for each $z_0 \in \Omega$ there is $r(z_0) > 0$ such that the disc 
$B(z_0, r(z_0)) = \{z : |z - z_0| < r(z_0)\}$ is contained in $\Omega$ and 
such that for every $r < r(z_0)$,
\begin{equation}\label{MVP of subharmonic functions}
  v(z_0) \leq \frac{1}{\pi r^2} \iint_{B(z_0,r)} v(z)\, dx\, dy.
\end{equation}
We often call (\ref{MVP of subharmonic functions}) the `\textbf{submean value property}'.
\end{itemize}
\end{definition}
\begin{remark}
Note that the integral in \ref{MVP of subharmonic functions} either
converges or diverges to $-\infty$, due to the upper-semicontinuity by the following
lemma:
\begin{lemma}\label{boundedness of uppersemicont fct}
Let $X$ be a topological space.
If $v:X\rightarrow [-\infty,\infty)$ is upper-semicontinuous, then $v$ is bounded
above on each compact subset $K\subset X$.
\end{lemma}
\end{remark}

Here are some examples of subharmonic functions:
\begin{example}
Every harmonic function is subharmonic, due to the \textbf{mean value property} of harmonic function:
\begin{theorem}
Let $\Omega$ be an open set in the plane and let $u$ be a harmonic function on $\Omega$. Then for
every $z_0\in\Omega$ and every $r>0$ such that the closed disc $\overline{B(z_0,r)}$ is contained
in $\Omega$,
\begin{equation*}
  u(z_0) = \frac{1}{2\pi}\int_0^{2\pi} u(z_0+re^{i\theta})\, d\theta
  = \frac{1}{\pi r^2} \iint_{B(z_0,r)} u(z)\, dx\, dy.
\end{equation*}
\end{theorem}
\end{example}
\begin{example}\label{log |f| is subharmonic}
Let $f(z)$ be an analytic function on a domain $\Omega\subset \mathbb{C}$, then 
$v(z)=\log |f(z)|$ is subharmonic on $\Omega$: 
\begin{itemize}
  \item $v$ is upper semicontinuous;
  \item For $z_0$ s.t. $f(z_0)=0$, condition \ref{MVP of subharmonic functions} trivially
  satisfied;
  \item For $z_0$ s.t. $f(z_0)\neq 0$, we can write $v(z)=\text{Re}(\log f(z))$ on some neighborhood
  of $z_0$, which is harmonic on this neighborhood.
\end{itemize}
\end{example}
\begin{example}
Let $f(z)$ be an analytic function on a domain $\Omega\subset \mathbb{C}$, then 
$|f(z)|^p$ is subharmonic on $\Omega$ as 
\begin{equation*}
  |f(z)|^p=\exp(p \log |f(z)|)
\end{equation*}
So the condition \ref{MVP of subharmonic functions} can be proved by the fact that
$\log |f(z)|$ is subharmonic and Jensen's inequality as $\exp$ is convex. More generally, we
have the following lemma:
\begin{lemma}
Given a subharmonic funciton $v$ on $\Omega$ and an increasing convex function $\varphi$ on
$[-\infty, \infty)$ which is continuous at $-\infty$, then $\varphi\circ v$ is subharmonic on
$\Omega$.
\end{lemma}
With this lemma, we can see that
\begin{equation*}
    \log^+ |f(z)|=\max(\log |f(z)|,0)
\end{equation*}
is also a subharmonic function.
\end{example}

There is an equivalent characterization of subharmonicity:
\begin{theorem}[Harmonic comparison characterization]\label{Harmonic comparison characterization}
Let $u:\Omega\rightarrow [-\infty,\infty)$ be an upper semicontinuous function. Then $u$ is 
subharmonic on $\Omega$ if and only if the following condition holds: If $v$ is a continuous function
on $\overline{B(z_0,r)}\subset \Omega$ that is harmonic on $B(z_0,r)$ s.t. for all 
$\xi\in \partial B(z_0,r)$, we have $u(\xi)\leq v(\xi)$, then 
\begin{equation*}
  u(z)\leq v(z),\quad z\in B(z_0,r).
\end{equation*}
\end{theorem}
\begin{proof}
(\textcolor{red}{missing})
\end{proof}

\begin{corollary}
If $\Omega$ is a connected open set and $v$ is a subharmonic function on $\Omega$ s.t. $v\not\equiv -\infty$,
then given any closed disc $\overline{B(z_0,r)}\subset \Omega$, 
\begin{equation*}
    \frac{1}{2\pi}\int v(z_0+r e^{i\theta})d\theta \neq -\infty.
\end{equation*}
\end{corollary}
\begin{proof}[Sketch of Proof]
First recall a lemma:
\begin{lemma}
Let $X$ be a metric space, $v$ be an upper-semicontinuous function. Then we can always find
a sequence $u_n$ of continuous function that decreasing to $v$ on $X$.
\end{lemma}
Now we pick $u_n$ to be a sequence of continuous functions decreasing to $v$ on 
$\partial B(z_0,r)$. Let $U_n(z)=P[u_n](z)$, i.e. the harmonic extension of $u_n$ to $B(z_0,r)$.
Assume for the sake of contradiction that we can find some $\overline{B(z_0,r)}\subset\Omega$ s.t.
\begin{equation*}
  \frac{1}{2\pi}\int v(z_0+r e^{i\theta})d\theta = -\infty.
\end{equation*}
Recall Lemma \ref{boundedness of uppersemicont fct}, $v$ is bounded above by some
constant $M\geq 0$. Write $v=v^+-v^-$ with $v^+=\max(v,0)$ and $v^-=\max(-v,0)$, so that
$0\leq v^+\leq M$ and $0\leq v^-\leq\infty$ on $\partial B(z_0,r)$. Since
\begin{equation*}
  \frac{1}{2\pi}\int v^+(z_0+re^{i\theta})\, d\theta\leq M<\infty,
\end{equation*}
the assumption is equivalent to
\begin{equation*}
  \frac{1}{2\pi}\int v^-(z_0+re^{i\theta})\, d\theta=\infty.
\end{equation*}
Now fix $z$ with $|z|<1$. Since $1-|z|\leq |e^{i\theta}-z|\leq 1+|z|$, we have
\begin{equation*}
  0<\frac{1-|z|}{1+|z|}\leq P_z(\theta)=\frac{1-|z|^2}{|e^{i\theta}-z|^2}\leq\frac{1+|z|}{1-|z|}
  \quad\text{for all }\theta.
\end{equation*}
Therefore
\begin{align*}
  \frac{1}{2\pi}\int P_z(\theta)\, v^+(z_0+re^{i\theta})\, d\theta
  &\leq \frac{1+|z|}{1-|z|}\, M<\infty,\\
  \frac{1}{2\pi}\int P_z(\theta)\, v^-(z_0+re^{i\theta})\, d\theta
  &\geq \frac{1-|z|}{1+|z|}\cdot\frac{1}{2\pi}\int v^-(z_0+re^{i\theta})\, d\theta=\infty.
\end{align*}
Therefore,
\begin{equation*}
  \frac{1}{2\pi}\int P_z(\theta)\, v(z_0+re^{i\theta})\, d\theta = -\infty,\quad |z|<1.
\end{equation*}
As
\begin{equation*}
  U_n(z)=\frac{1}{2\pi}\int P_z(\theta)\, u_n(z_0+re^{i\theta})\, d\theta 
\end{equation*}
and $u_n\downarrow v$, so $U_n(z)\rightarrow -\infty$ for each $z\in B(z_0,r)$. By Harmonic
comparison \ref{Harmonic comparison characterization}, we know that $U_n \geq v$ over $B(z_0,r)$,
so $v=-\infty$ on $B(z_0,r)$. With above result, we can actually extend this local property globally
using standard topological argument as follow: consider 
\begin{equation*}
  A:=\{z\in \Omega: v(z)\equiv -\infty\text{ on a neighborhood of }z\}
\end{equation*}
then one can show that it is open and closed. As $z_0\in A$, it is nonempty, thus $v\equiv -\infty$
over $\Omega$, which gives a contradiction.

\end{proof}

To investigate further properties of subharmonic functions, we'd like to first define
an auxiliary function called \textbf{Poisson modification}:
\begin{definition}[Poisson modification]\label{Poisson modification}
Let $v:\mathbb{D}\rightarrow[-\infty,\infty)$ be an upper semicontinuous function and let
$0<r<1$. The \textbf{Poisson modification} of $v$ on $B(0,r)$ is the function
\begin{equation*}
  v_r(z):=
  \begin{cases}
    v(z), & r\leq|z|<1,\\[4pt]
    \displaystyle\frac{1}{2\pi}\int P_{z/r}(\theta)\, v(re^{i\theta})\, d\theta, & |z|<r.
  \end{cases}
\end{equation*}
\end{definition}
By Lemma \ref{boundedness of uppersemicont fct}, $v$ is bounded above on $\partial B(0,r)$, so
the integral is well-defined with values in $[-\infty,\infty)$.

\begin{theorem}\label{Properties of Poisson modification}
Let $v$ be a subharmonic function on $\mathbb{D}$ s.t. $v\not\equiv-\infty$, and let $0<r<1$.
Then
\begin{itemize}
  \item[(a)] $v_r$ is finite and harmonic on $B(0,r)$;
  \item[(b)] $v_r$ is subharmonic on $\mathbb{D}$;
  \item[(c)] $v(z)\leq v_r(z)$ for all $z\in\mathbb{D}$;
  \item[(d)] $v_r(z)$ is an increasing function of $r$.
\end{itemize}
\end{theorem}
\begin{proof}
(\textcolor{red}{missing})
\end{proof}
If we take $z=0$ in Theorem \ref{Properties of Poisson modification}(d), we get:
\begin{corollary}
If $v(z)$ is a subharmonic function on $\mathbb{D}$, then
\begin{equation*}
  m(r)=v_r(0)=\frac{1}{2\pi}\int v(r e^{i\theta})d\theta
\end{equation*}
is an increasing function of $r\in(0,1)$.
\end{corollary}

\begin{definition}[Harmonic majorant]\label{Harmonic majorant}
Let $v$ be a subharmonic function on an open set $\Omega$. A \textbf{harmonic majorant} of $v$
is a harmonic function $U$ on $\Omega$ s.t. $v(z)\leq U(z)$ for all $z\in\Omega$. A harmonic
majorant $u$ of $v$ is called the \textbf{least harmonic majorant} of $v$ if $u(z)\leq U(z)$,
$z\in\Omega$, for every harmonic majorant $U$ of $v$.
\end{definition}

\begin{theorem}\label{Harmonic majorant criterion}
Let $v$ be a subharmonic function on $\mathbb{D}$ s.t. $v\not\equiv-\infty$. Then $v$ has a
harmonic majorant if and only if
\begin{equation*}
  \sup_{0<r<1}\frac{1}{2\pi}\int v(re^{i\theta})\, d\theta=\sup_{0<r<1}v_r(0)<\infty.
\end{equation*}
In this case, the least harmonic majorant of $v$ is
\begin{equation*}
  u(z)=\lim_{r\rightarrow 1}\frac{1}{2\pi}\int P_{z/r}(\theta)\, v(re^{i\theta})\, d\theta
  =\lim_{r\rightarrow 1}v_r(z),\quad z\in\mathbb{D}.
\end{equation*}
\end{theorem}

\begin{proof}
\textcolor{red}{missing}
\end{proof}.

\chapter{Background}
\label{sec:background}
Before we talk about the main object in this thesis, we need to first introduce some analysis
backgrounds.
\section{Poisson integral}
Let $\mathbb{D}:=\{z\in\C: |z|<1\}$ and $\mathbb{T}=\partial \mathbb{D}$. 
When given a continuous function $f$ on $\mathbb{T}$, the famous \textbf{Dirichlet
problem}, which is to find the solution for the following equations:
\begin{equation*}
  \begin{cases}
  \Delta u=0, & z\in \mathbb{D}\\
  u=f, & z\in \mathbb{T}.
  \end{cases}
\end{equation*}
is asking for a harmonic extension $u$ over $\mathbb{D}$ with the boundary data $f$ on $\mathbb{T}$.
To give a solution for this problem, we'd like to introduce the notion called \textbf{Poisson integral}.
This section is based on \cite{garnett2006bounded}, I.3.
\begin{definition}
Let $c\in\mathbb{C}$. The \textbf{Herglotz--Riesz kernel centered at $c$} is the function
\begin{equation*}
  K(c;z,w) := \frac{(w-c)+(z-c)}{(w-c)-(z-c)} = \frac{w+z-2c}{w-z},
  \qquad z,w\in\mathbb{C},\ w\neq z.
\end{equation*}
The \textbf{Poisson kernel centered at $c$} is defined as
\begin{equation*}
  P(c;z,w) := \text{Re}\; K(c;z,w)= \frac{|w-c|^2-|z-c|^2}{|w-z|^2},
  \qquad z,w\in\mathbb{C},\ w\neq z.
\end{equation*}
and we denote $K(z,w):= K(0;z,w)$, $P(z,w):= P(0;z,w)$, and $K_z(\theta):= K(z,e^{i\theta})$,
 $P_z(\theta):= P(z,e^{i\theta})$. 
\end{definition}

\begin{definition}
Given a function $f(\theta)\in L^1(\mathbb{T})$, we define the \textbf{Poisson integral} of $f$ 
to be the function
\begin{equation*}
  P[f](z):=\frac{1}{2\pi}\int P_z(\theta) f(\theta)d\theta.
\end{equation*}
In general, given a finite (signed) measure $\mu$ on $\mathbb{T}$, we define the \textbf{Poisson integral} of $\mu$
to be the function
\begin{equation*}
  P[\mu](z):=\int P_z(\theta) d\mu (\theta).
\end{equation*}
The first definition is a special case of the second one if we set 
$d\mu(\theta)=\frac{1}{2\pi} f(\theta)d\theta$.
\end{definition}

\begin{theorem}[Poisson integral formula]\label{Poisson integral formula}
Let $u$ be a continuous function on the closed unit disc $\overline{\mathbb{D}}$. If $u$ is
harmonic on $\mathbb{D}$, then
\begin{equation*}
  u(z)=\frac{1}{2\pi}\int_0^{2\pi} P_z(\theta)\, u(e^{i\theta})\, d\theta
  \qquad \text{for all } z\in\mathbb{D}.
\end{equation*}
\end{theorem}
\begin{proof}
Recall that if $u$ is harmonic on $\mathbb{D}$, then it has the 
\textbf{mean value property}: 
\begin{equation*}
  u(0)=\frac{1}{2\pi}\int_0^{2\pi} u(e^{i\theta})d\theta
\end{equation*}
Now given $z_0=re^{i\theta_0}$, consider the M\"{o}bius transformation 
$\tau(z)=(z-z_0)/(1-\bar{z}_0 z)$, and write $\tau(e^{i\theta})=e^{i\varphi}$.
Then we can compute that $d\varphi/d\theta=P_{z_0}(\theta)$, therefore
\begin{equation*}
  u(z_0)=u(\tau^{-1}(0))=\frac{1}{2\pi}\int_0^{2\pi} (u\circ\tau^{-1})(e^{i\varphi})
  d\varphi=\frac{1}{2\pi}\int_0^{2\pi} u(e^{i\theta})P_{z_0}(\theta)d\theta.
\end{equation*}
\end{proof}
Moreover, we have part of the converse of above theorem:
\begin{theorem}\label{Poisson integral of L^1 is harmonic}
Given a function $f(\theta)\in L^1(\mathbb{T})$, its Poisson integral $P[f]$
is harmonic on $\mathbb{D}$.
\end{theorem}
\begin{proof}[Sketch of Proof]
Recall that a function $u(z)$ is said to be harmonic if 
\begin{equation*}
  \Delta u=\frac{\partial^2 u}{\partial x^2}+\frac{\partial^2 u}{\partial y^2}=0.
\end{equation*}
As Herglotz--Riesz kernel $K_z(\theta)$ is analytic over $\mathbb{D}$, 
the Poisson kernel $P_z(\theta)=\text{Re} K_z(\theta)$ is harmonic over 
$\mathbb{D}$, thus for $z\in \mathbb{D}$,
\begin{equation*}
  \Delta P[f](z)=\frac{1}{2\pi}\int \cancel{\Delta P_z(\theta)} f(\theta)d\theta=0.
\end{equation*}
\end{proof}
More generally, the same holds for Poisson integrals of measures:
\begin{theorem}
Given a finite measure $\mu$ on $\mathbb{T}$, its Poisson integral $P[\mu]$
is harmonic on $\mathbb{D}$.
\end{theorem}
\begin{proof}[Sketch of Proof]
Similar as in the previous proof, for $z\in\mathbb{D}$,
\begin{equation*}
  \Delta P[\mu](z)=\int \cancel{\Delta P_z(\theta)}\, d\mu(\theta)=0.
\end{equation*}
\end{proof}
The full converse theorem of Poisson integral formula states that the Poisson integral indeed
gives a continuous solution for Dirichlet problem, which will be later proved in Corollary 
\ref{Dirichlet problem on the disc}. Before that, we'd like to introduce more properties of 
Poisson integral.



Here are some properties for Poisson kernel:
\begin{lemma}\label{Properties of Poisson kernel}
Let $z\in\mathbb{D}$. Then
\begin{itemize}
  \item[(a)] $P_z(\theta)>0$ for all $\theta$, and
  $\displaystyle\frac{1}{2\pi}\int_0^{2\pi} P_z(\theta)\, d\theta=1$;
  \item[(b)] for any $\theta_0\in\mathbb{R}$ and $0<\delta<\pi$,
  \begin{equation*}
    \sup_{\delta\leq|\theta-\theta_0|\leq\pi}P_z(\theta)\rightarrow 0
    \quad\text{and}\quad
    \frac{1}{2\pi}\int_{\delta\leq|\theta-\theta_0|\leq\pi}P_z(\theta)\, d\theta\rightarrow 0
    \qquad (z\rightarrow e^{i\theta_0}).
  \end{equation*}
\end{itemize}
\end{lemma}

Here's an important estimate for Poisson integral of $L^p$ function:
\begin{theorem}
Let $1\leq p\leq\infty$ and let $f$ be a function in $L^p(\mathbb{T})$, and
$\|f\|_p^p:=\frac{1}{2\pi}\int_0^{2\pi}|f(\theta)|^p\, d\theta$ for $p<\infty$. Then for every
$0<r<1$,
\begin{equation*}
  \|P[f](re^{i\cdot})\|_p\leq\|f\|_p,
\end{equation*}
and hence $\sup_{0<r<1}\|P[f](re^{i\cdot})\|_p\leq\|f\|_p$.
\end{theorem}
\begin{proof}[Sketch of Proof]
If $p=\infty$, then $|P[f](z)|\leq\frac{1}{2\pi}\int P_z(\theta)|f(\theta)|\, d\theta
\leq\|f\|_\infty$.

If $1\leq p<\infty$, for $z\in\mathbb{D}$, part (a) of the lemma \ref{Properties of Poisson kernel} says that $P_z(\theta)\frac{d\theta}{2\pi}$ is a
probability measure on $\mathbb{T}$. Jensen's inequality for this probability measure gives
\begin{equation*}
  |P[f](z)|^p\leq\frac{1}{2\pi}\int_0^{2\pi} P_z(\theta)|f(\theta)|^p\, d\theta.
\end{equation*}
Set $z=re^{i\varphi}$. Note the symmetry of Poisson kernel:
\begin{equation*}
  P_{re^{i\varphi}}(\theta)=\frac{1-r^2}{1-2r\cos(\theta-\varphi)+r^2}=P_{re^{i\theta}}(\varphi),
\end{equation*}
Thus:
\begin{equation*}
  \frac{1}{2\pi}\int_0^{2\pi}|P[f](re^{i\varphi})|^p\, d\varphi
  \leq\frac{1}{2\pi}\int_0^{2\pi}|f(\theta)|^p
  \left(\frac{1}{2\pi}\int_0^{2\pi}P_{re^{i\theta}}(\varphi)\, d\varphi\right) d\theta
  =\|f\|_p^p. \qedhere
\end{equation*}
\end{proof}

The following theorem shows that the Poisson kernel acts like the identity as $r\rightarrow 1$:
\begin{theorem}\label{Poisson approximate identity}
Let $C(\mathbb{T})$ be the space of all continuous functions over $\mathbb{T}$. Then
\begin{itemize}
  \item[(a)] If $f\in C(\mathbb{T})$, then $P[f](re^{i\theta})$ converges uniformly to $f(e^{i\theta})$
  as $r\rightarrow 1$.
  \item[(b)] If $\mu$ is a finite measure on $\mathbb{T}$, then the measures
  $P[\mu](re^{i\theta})\, d\theta/(2\pi)$ converge weak-star to $\mu$ as $r\rightarrow 1$, i.e.
  \begin{equation*}
    \frac{1}{2\pi}\int_0^{2\pi} g(\theta)P[\mu](re^{i\theta})\, d\theta\rightarrow\int g\, d\mu
    \qquad (r\rightarrow 1)
  \end{equation*}
  for every $g\in C(\mathbb{T})$.
\end{itemize}
\end{theorem}
\begin{proof}[Sketch of Proof]
(a) Fix $\varphi$. By part (a) of the lemma \ref{Properties of Poisson kernel},
\begin{equation*}
  |P[f](re^{i\varphi})-f(e^{i\varphi})|
  \leq\frac{1}{2\pi}\int_{|\theta-\varphi|<\delta}P_{re^{i\varphi}}(\theta)|f(e^{i\theta})-f(e^{i\varphi})|\, d\theta
  +\frac{1}{2\pi}\int_{\delta\leq|\theta-\varphi|\leq\pi}P_{re^{i\varphi}}(\theta)|f(e^{i\theta})-f(e^{i\varphi})|\, d\theta.
\end{equation*}
As $\mathbb{T}$ is compact, $f$ is uniformly continuous on $\mathbb{T}$.
Given $\varepsilon>0$, by uniform continuity of $f$, we can pick $\delta>0$ s.t.
$|f(e^{i\theta})-f(e^{i\varphi})|<\varepsilon$ for all $\varphi$ and all $|\theta-\varphi|<\delta$,
then we have
\begin{equation*}
  \frac{1}{2\pi}\int_{|\theta-\varphi|<\delta}P_{re^{i\varphi}}(\theta)|f(e^{i\theta})-f(e^{i\varphi})|\, d\theta
  \leq \varepsilon\cdot\frac{1}{2\pi}\int P_{re^{i\varphi}}(\theta)d\theta=\varepsilon.
\end{equation*}
For this $\delta$, the
second term is at most $2\|f\|_\infty\cdot\frac{1}{2\pi}\int_{\delta\leq|\theta-\varphi|\leq\pi}
P_{re^{i\varphi}}(\theta)\, d\theta$. Since $P_{re^{i\varphi}}(\theta)=P_r(\theta-\varphi)$, this integral
does not depend on $\varphi$, and it tends to $0$ as $r\rightarrow 1$ by part (b) of the
lemma \ref{Properties of Poisson kernel}. Hence the convergence is uniform in $\varphi$.

(b) Let $g\in C(\mathbb{T})$. For fixed $r$, the Poisson kernel is bounded, so by Fubini's theorem and
the symmetry $P_{re^{i\theta}}(\varphi)=P_{re^{i\varphi}}(\theta)$,
\begin{equation*}
  \frac{1}{2\pi}\int_0^{2\pi} g(\theta)P[\mu](re^{i\theta})\, d\theta
  =\int\left(\frac{1}{2\pi}\int_0^{2\pi}P_{re^{i\varphi}}(\theta)g(\theta)\, d\theta\right) d\mu(\varphi)
  =\int P[g](re^{i\varphi})\, d\mu(\varphi).
\end{equation*}
By (a), $P[g](re^{i\varphi})\rightarrow g(e^{i\varphi})$ uniformly, and $\mu$ is finite, so the
right-hand side tends to $\int g\, d\mu$.
\end{proof}

As a consequence, the Poisson integral solves the Dirichlet problem on the disc (which is 
the converse result of the Poisson integral formula
 \ref{Poisson integral formula}):
\begin{corollary}[Dirichlet problem on the disc]\label{Dirichlet problem on the disc}
Let $f$ be a continuous function on $\mathbb{T}$. Define
\begin{equation*}
  u(z):=
  \begin{cases}
    P[f](z), & z\in\mathbb{D},\\
    f(z), & z\in\mathbb{T}.
  \end{cases}
\end{equation*}
Then $u$ is continuous on $\overline{\mathbb{D}}$ and harmonic on $\mathbb{D}$, i.e. 
a continuous solution on $\overline{\mathbb{D}}$ for the Dirichlet
problem:
\begin{equation*}
  \begin{cases}
  \Delta u=0, & z\in \mathbb{D}\\
  u=f, & z\in \mathbb{T}.
  \end{cases}
\end{equation*}
\end{corollary}
\begin{proof}
$u$ is harmonic on $\mathbb{D}$ by Theorem \ref{Poisson integral of L^1 is harmonic}. For
continuity at $e^{i\theta_0}\in\mathbb{T}$, let $z=re^{i\varphi}\in\mathbb{D}$; then
\begin{equation*}
  |u(z)-f(e^{i\theta_0})|\leq|P[f](re^{i\varphi})-f(e^{i\varphi})|+|f(e^{i\varphi})-f(e^{i\theta_0})|,
\end{equation*}
which tends to $0$ as $z\rightarrow e^{i\theta_0}$ by part (a) of
Theorem \ref{Poisson approximate identity} and the continuity of $f$.
\end{proof}

Moreover, a measure on $\mathbb{T}$ is determined by its Poisson integral:
\begin{corollary}
Let $\mu$ and $\nu$ be finite measures on $\mathbb{T}$. If $P[\mu]=P[\nu]$ on $\mathbb{D}$,
then $\mu=\nu$.
\end{corollary}
\begin{proof}
Let $\lambda:=\mu-\nu$. By linearity, $P[\lambda]=0$ on $\mathbb{D}$. For every $g\in C(\mathbb{T})$,
part (b) of Theorem \ref{Poisson approximate identity} gives
\begin{equation*}
  \int g\, d\lambda=\lim_{r\rightarrow 1}\frac{1}{2\pi}\int_0^{2\pi} g(\theta)P[\lambda](re^{i\theta})\, d\theta=0.
\end{equation*}
By the uniqueness part of the Riesz representation theorem, $\lambda=0$, thus $\mu=\nu$.
\end{proof}

Finally, Poisson integrals of $L^p$ functions and of measures are characterized by the growth
of their radial means:
\begin{theorem}\label{Characterization of Poisson integrals}
Let $u$ be a harmonic function on $\mathbb{D}$.
\begin{itemize}
  \item[(a)] If $1<p\leq\infty$, then $u=P[f]$ for some $f\in L^p(\mathbb{T})$ if and only if
  \begin{equation*}
    \sup_{0<r<1}\|u(re^{i\cdot})\|_p<\infty.
  \end{equation*}
  \item[(b)] $u=P[\mu]$ for some finite measure $\mu$ on $\mathbb{T}$ if and only if
  \begin{equation*}
    \sup_{0<r<1}\|u(re^{i\cdot})\|_1<\infty.
  \end{equation*}
  \item[(c)] (Herglotz) $u\geq 0$ on $\mathbb{D}$ if and only if $u=P[\mu]$ for some finite
  positive measure $\mu$ on $\mathbb{T}$.
\end{itemize}
\end{theorem}

\begin{proof}
(\textcolor{red}{missing})
\end{proof}

\section{Nontangential limit}
When we are discussing functions over $\mathbb{D}$, it will be good to have the proper
notion of `limit function' over the boundary $\mathbb{T}=\partial \mathbb{D}$.
This section is mainly based on \cite{rudin1974real}, 11.21 - 11.24.
\begin{definition}
Fix $\alpha>1$. For $\zeta\in\mathbb{T}$, we consider the \textbf{Stolz region}:
\begin{equation*}
  \Gamma_\alpha(\zeta):=\left\{z\in\mathbb{D}: |\zeta-z|<\alpha(1-|z|)\right\}.
\end{equation*}
A function $u:\mathbb{D}\to\mathbb{C}$ is said to have \textbf{nontangential limit} $\lambda\in\mathbb{C}$
at $\zeta\in\mathbb{T}$ if, for every $\alpha>1$,
\begin{equation*}
  \lim_{\substack{z\to\zeta\\ z\in\Gamma_\alpha(\zeta)}}u(z)=L.
\end{equation*}
When this limit exists, we write $u^*(\zeta)=\lambda$ and call it the \textbf{(nontangential) boundary value} of $u$ at $\zeta$.
This give rise to a function $u^*:\mathbb{T}\rightarrow\mathbb{C}$, for which we call the \textbf{(nontangential) boundary function}
of $u$.
\end{definition}


Before we discuss the main results related with nontangential limit, we first 
introduce a notion called \textbf{Lebesgue point}:
\begin{definition}[Lebesgue point]
Consider $(\mathbb{T},d\sigma)$, where $d\sigma := d\theta/(2\pi)$.
If $f\in L^1(\mathbb{T})$, a point $\zeta=e^{i\theta_0}\in\mathbb{T}$
is called a \textbf{Lebesgue point} of $f$ if
\begin{equation*}
  \lim_{h\downarrow 0}
  \frac{1}{2h}
  \int_{\theta_0-h}^{\theta_0+h}
  \left|f(e^{i\theta})-f(e^{i\theta_0})\right|\,d\theta=0.
\end{equation*}
\end{definition}

\begin{remark}
If $f$ is continuous at $\xi$, then $\xi$ is a Lebesgue point of $f$. So one can
think of Lebesgue points of a function are those points where in average sense, the
function does not oscillate too much.
\end{remark}

\begin{theorem}[Lebesgue differentiation theorem]
\label{thm:lebesgue_differentiation}
Let $f\in L^1(\mathbb{T})$. Then almost every point of $\mathbb{T}$
 is a Lebesgue point of $f$.
\end{theorem}

The following result tells us that nontangential limit of
the Poisson integral of a function gives back the function value itself
at all lebesgue points of that function:
\begin{theorem}
\label{thm:poisson_nontangential_lebesgue_point}
Let $f\in L^1(\mathbb{T})$ be real-valued. If $\zeta\in\mathbb{T}$ is a
Lebesgue point of $f$, then the Poisson integral $P[f]$ has nontangential
limit $f(\zeta)$ at $\zeta$. That is, for every $\alpha>1$,
\begin{equation*}
  \lim_{\substack{z\to\zeta\\ z\in\Gamma_\alpha(\zeta)}}P[f](z)=f(\zeta).
\end{equation*}
\end{theorem}

Then using Lebesgue differentiation theorem \ref{thm:lebesgue_differentiation},
we get an a.e. recovery of boundary data:
\begin{corollary}[Almost everywhere recovery of boundary data]
\label{cor:poisson_nontangential_ae}
Let $f\in L^1(\mathbb{T})$ be real-valued. Then $P[f]$ has nontangential
limits almost everywhere on $\mathbb{T}$, and
\begin{equation*}
  (P[f])^*(\zeta)=f(\zeta)
  \qquad\text{for almost every }\zeta\in\mathbb{T}.
\end{equation*}
\end{corollary}

\begin{theorem}
\label{thm:singular_poisson_nontangential_zero}
Let $\nu$ be a finite measure on $\mathbb{T}$ s.t. $\nu\bot d\sigma$, i.e. singular
with respect to $d\sigma = d\theta/(2\pi)$. Then its Poisson integral $P[\nu]$ has
nontangential limit 0 at almost every point of $\mathbb{T}$. For every $\alpha>1$,
\begin{equation*}
  \lim_{\substack{z\to\zeta\\ z\in\Gamma_\alpha(\zeta)}}P[\nu](z)=0
  \qquad\text{for a.e. }\zeta\in\mathbb{T}.
\end{equation*}
\end{theorem}

\begin{theorem}[Fatou's theorem for finite measures]
\label{thm:fatou_signed_measure}
Let $\mu$ be a finite measure on $\mathbb{T}$ with Lebesgue
decomposition
\begin{equation*}
  d\mu=f\,\frac{d\theta}{2\pi}+d\mu_s,
  \qquad f\in L^1(\mathbb{T}),
  \qquad \mu_s\perp\frac{d\theta}{2\pi}.
\end{equation*}
Then $P[\mu]$ has nontangential limits a.e. on $\mathbb{T}$, and for every $\alpha>1$,
\begin{equation*}
  \lim_{\substack{z\to\zeta\\ z\in\Gamma_\alpha(\zeta)}}P[\mu](z)=f(\zeta)
  \qquad\text{for a.e. }\zeta\in\mathbb{T}.
\end{equation*}
\end{theorem}
\begin{proof}
By linearity, $P[\mu]=P[f]+P[\mu_s]$. Apply
Corollary~\ref{cor:poisson_nontangential_ae} and
Theorem~\ref{thm:singular_poisson_nontangential_zero} will get the desired result.
\end{proof}





\section{Subharmonic function}
\label{subsec:subharmonic}
To investigate the main objects in this thesis,
which are Hardy spaces and Nevanlinna class, it will be good to know the properties of 
function like $|f|^p$, $\log |f|$ and $\log^+ |f|$
for analytic function $f$. These functions belong to an important class of functions called 
\textbf{subharmonic functions}. This section is based on \cite{garnett2006bounded}, I.6.
\begin{definition}
Let $\Omega$ be an open set in the plane. A \textbf{subharmonic function} on $\Omega$ is a function
$v : \Omega \to [-\infty, \infty)$ such that
\begin{itemize}
  \item[(a)] $v$ is \emph{upper semicontinous}, i.e. $v(z_0)\geq \lim_{z\rightarrow z_0}v(z)$ 
  for all $z_0\in\Omega$;
  \item[(b)] for each $z_0 \in \Omega$ there is $r(z_0) > 0$ such that the disc 
$B(z_0, r(z_0)) = \{z : |z - z_0| < r(z_0)\}$ is contained in $\Omega$ and 
such that for every $r < r(z_0)$,
\begin{equation}\label{MVP of subharmonic functions}
  v(z_0) \leq \frac{1}{\pi r^2} \iint_{B(z_0,r)} v(z)\, dx\, dy.
\end{equation}
We often call (\ref{MVP of subharmonic functions}) the `\textbf{submean value property}'.
\end{itemize}
\end{definition}
\begin{remark}
Note that the integral in \ref{MVP of subharmonic functions} either
converges or diverges to $-\infty$, due to the upper-semicontinuity by the following
lemma:
\begin{lemma}\label{boundedness of uppersemicont fct}
Let $X$ be a topological space.
If $v:X\rightarrow [-\infty,\infty)$ is upper-semicontinuous, then $v$ is bounded
above on each compact subset $K\subset X$.
\end{lemma}
\end{remark}

Here are some examples of subharmonic functions:
\begin{example}
Every harmonic function is subharmonic, due to the \textbf{mean value property} of harmonic function:
\begin{theorem}
Let $\Omega$ be an open set in the plane and let $u$ be a harmonic function on $\Omega$. Then for
every $z_0\in\Omega$ and every $r>0$ such that the closed disc $\overline{B(z_0,r)}$ is contained
in $\Omega$,
\begin{equation*}
  u(z_0) = \frac{1}{2\pi}\int_0^{2\pi} u(z_0+re^{i\theta})\, d\theta
  = \frac{1}{\pi r^2} \iint_{B(z_0,r)} u(z)\, dx\, dy.
\end{equation*}
\end{theorem}
\end{example}
\begin{example}\label{log |f| is subharmonic}
Let $f(z)$ be an analytic function on a domain $\Omega\subset \mathbb{C}$, then 
$v(z)=\log |f(z)|$ is subharmonic on $\Omega$: 
\begin{itemize}
  \item $v$ is upper semicontinuous;
  \item For $z_0$ s.t. $f(z_0)=0$, condition \ref{MVP of subharmonic functions} trivially
  satisfied;
  \item For $z_0$ s.t. $f(z_0)\neq 0$, we can write $v(z)=\text{Re}(\log f(z))$ on some neighborhood
  of $z_0$, which is harmonic on this neighborhood.
\end{itemize}
\end{example}
\begin{example}
Let $f(z)$ be an analytic function on a domain $\Omega\subset \mathbb{C}$, then 
$|f(z)|^p$ is subharmonic on $\Omega$ as 
\begin{equation*}
  |f(z)|^p=\exp(p \log |f(z)|)
\end{equation*}
So the condition \ref{MVP of subharmonic functions} can be proved by the fact that
$\log |f(z)|$ is subharmonic and Jensen's inequality as $\exp$ is convex. More generally, we
have the following lemma:
\begin{lemma}
Given a subharmonic funciton $v$ on $\Omega$ and an increasing convex function $\varphi$ on
$[-\infty, \infty)$ which is continuous at $-\infty$, then $\varphi\circ v$ is subharmonic on
$\Omega$.
\end{lemma}
With this lemma, we can see that
\begin{equation*}
    \log^+ |f(z)|=\max(\log |f(z)|,0)
\end{equation*}
is also a subharmonic function.
\end{example}

There is an equivalent characterization of subharmonicity:
\begin{theorem}[Harmonic comparison characterization]\label{Harmonic comparison characterization}
Let $u:\Omega\rightarrow [-\infty,\infty)$ be an upper semicontinuous function. Then $u$ is 
subharmonic on $\Omega$ if and only if the following condition holds: If $v$ is a continuous function
on $\overline{B(z_0,r)}\subset \Omega$ that is harmonic on $B(z_0,r)$ s.t. for all 
$\xi\in \partial B(z_0,r)$, we have $u(\xi)\leq v(\xi)$, then 
\begin{equation*}
  u(z)\leq v(z),\quad z\in B(z_0,r).
\end{equation*}
\end{theorem}
\begin{proof}
(\textcolor{red}{missing})
\end{proof}

\begin{corollary}
If $\Omega$ is a connected open set and $v$ is a subharmonic function on $\Omega$ s.t. $v\not\equiv -\infty$,
then given any closed disc $\overline{B(z_0,r)}\subset \Omega$, 
\begin{equation*}
    \frac{1}{2\pi}\int v(z_0+r e^{i\theta})d\theta \neq -\infty.
\end{equation*}
\end{corollary}
\begin{proof}[Sketch of Proof]
First recall a lemma:
\begin{lemma}
Let $X$ be a metric space, $v$ be an upper-semicontinuous function. Then we can always find
a sequence $u_n$ of continuous function that decreasing to $v$ on $X$.
\end{lemma}
Now we pick $u_n$ to be a sequence of continuous functions decreasing to $v$ on 
$\partial B(z_0,r)$. Let $U_n(z)=P[u_n](z)$, i.e. the harmonic extension of $u_n$ to $B(z_0,r)$.
Assume for the sake of contradiction that we can find some $\overline{B(z_0,r)}\subset\Omega$ s.t.
\begin{equation*}
  \frac{1}{2\pi}\int v(z_0+r e^{i\theta})d\theta = -\infty.
\end{equation*}
Recall Lemma \ref{boundedness of uppersemicont fct}, $v$ is bounded above by some
constant $M\geq 0$. Write $v=v^+-v^-$ with $v^+=\max(v,0)$ and $v^-=\max(-v,0)$, so that
$0\leq v^+\leq M$ and $0\leq v^-\leq\infty$ on $\partial B(z_0,r)$. Since
\begin{equation*}
  \frac{1}{2\pi}\int v^+(z_0+re^{i\theta})\, d\theta\leq M<\infty,
\end{equation*}
the assumption is equivalent to
\begin{equation*}
  \frac{1}{2\pi}\int v^-(z_0+re^{i\theta})\, d\theta=\infty.
\end{equation*}
Now fix $z$ with $|z|<1$. Since $1-|z|\leq |e^{i\theta}-z|\leq 1+|z|$, we have
\begin{equation*}
  0<\frac{1-|z|}{1+|z|}\leq P_z(\theta)=\frac{1-|z|^2}{|e^{i\theta}-z|^2}\leq\frac{1+|z|}{1-|z|}
  \quad\text{for all }\theta.
\end{equation*}
Therefore
\begin{align*}
  \frac{1}{2\pi}\int P_z(\theta)\, v^+(z_0+re^{i\theta})\, d\theta
  &\leq \frac{1+|z|}{1-|z|}\, M<\infty,\\
  \frac{1}{2\pi}\int P_z(\theta)\, v^-(z_0+re^{i\theta})\, d\theta
  &\geq \frac{1-|z|}{1+|z|}\cdot\frac{1}{2\pi}\int v^-(z_0+re^{i\theta})\, d\theta=\infty.
\end{align*}
Therefore,
\begin{equation*}
  \frac{1}{2\pi}\int P_z(\theta)\, v(z_0+re^{i\theta})\, d\theta = -\infty,\quad |z|<1.
\end{equation*}
As
\begin{equation*}
  U_n(z)=\frac{1}{2\pi}\int P_z(\theta)\, u_n(z_0+re^{i\theta})\, d\theta 
\end{equation*}
and $u_n\downarrow v$, so $U_n(z)\rightarrow -\infty$ for each $z\in B(z_0,r)$. By Harmonic
comparison \ref{Harmonic comparison characterization}, we know that $U_n \geq v$ over $B(z_0,r)$,
so $v=-\infty$ on $B(z_0,r)$. With above result, we can actually extend this local property globally
using standard topological argument as follow: consider 
\begin{equation*}
  A:=\{z\in \Omega: v(z)\equiv -\infty\text{ on a neighborhood of }z\}
\end{equation*}
then one can show that it is open and closed. As $z_0\in A$, it is nonempty, thus $v\equiv -\infty$
over $\Omega$, which gives a contradiction.

\end{proof}

To investigate further properties of subharmonic functions, we'd like to first define
an auxiliary function called \textbf{Poisson modification}:
\begin{definition}[Poisson modification]\label{Poisson modification}
Let $v:\mathbb{D}\rightarrow[-\infty,\infty)$ be an upper semicontinuous function and let
$0<r<1$. The \textbf{Poisson modification} of $v$ on $B(0,r)$ is the function
\begin{equation*}
  v_r(z):=
  \begin{cases}
    v(z), & r\leq|z|<1,\\[4pt]
    \displaystyle\frac{1}{2\pi}\int P_{z/r}(\theta)\, v(re^{i\theta})\, d\theta, & |z|<r.
  \end{cases}
\end{equation*}
\end{definition}
By Lemma \ref{boundedness of uppersemicont fct}, $v$ is bounded above on $\partial B(0,r)$, so
the integral is well-defined with values in $[-\infty,\infty)$.

\begin{theorem}\label{Properties of Poisson modification}
Let $v$ be a subharmonic function on $\mathbb{D}$ s.t. $v\not\equiv-\infty$, and let $0<r<1$.
Then
\begin{itemize}
  \item[(a)] $v_r$ is finite and harmonic on $B(0,r)$;
  \item[(b)] $v_r$ is subharmonic on $\mathbb{D}$;
  \item[(c)] $v(z)\leq v_r(z)$ for all $z\in\mathbb{D}$;
  \item[(d)] $v_r(z)$ is an increasing function of $r$.
\end{itemize}
\end{theorem}
\begin{proof}
(\textcolor{red}{missing})
\end{proof}
If we take $z=0$ in Theorem \ref{Properties of Poisson modification}(d), we get:
\begin{corollary}
If $v(z)$ is a subharmonic function on $\mathbb{D}$, then
\begin{equation*}
  m(r)=v_r(0)=\frac{1}{2\pi}\int v(r e^{i\theta})d\theta
\end{equation*}
is an increasing function of $r\in(0,1)$.
\end{corollary}

\begin{definition}[Harmonic majorant]\label{Harmonic majorant}
Let $v$ be a subharmonic function on an open set $\Omega$. A \textbf{harmonic majorant} of $v$
is a harmonic function $U$ on $\Omega$ s.t. $v(z)\leq U(z)$ for all $z\in\Omega$. A harmonic
majorant $u$ of $v$ is called the \textbf{least harmonic majorant} of $v$ if $u(z)\leq U(z)$,
$z\in\Omega$, for every harmonic majorant $U$ of $v$.
\end{definition}

\begin{theorem}\label{Harmonic majorant criterion}
Let $v$ be a subharmonic function on $\mathbb{D}$ s.t. $v\not\equiv-\infty$. Then $v$ has a
harmonic majorant if and only if
\begin{equation*}
  \sup_{0<r<1}\frac{1}{2\pi}\int v(re^{i\theta})\, d\theta=\sup_{0<r<1}v_r(0)<\infty.
\end{equation*}
In this case, the least harmonic majorant of $v$ is
\begin{equation*}
  u(z)=\lim_{r\rightarrow 1}\frac{1}{2\pi}\int P_{z/r}(\theta)\, v(re^{i\theta})\, d\theta
  =\lim_{r\rightarrow 1}v_r(z),\quad z\in\mathbb{D}.
\end{equation*}
\end{theorem}

\begin{proof}
\textcolor{red}{missing}
\end{proof}.

\chapter{Background}
\label{sec:background}
Before we talk about the main object in this thesis, we need to first introduce some analysis
backgrounds.
\section{Poisson integral}
Let $\mathbb{D}:=\{z\in\C: |z|<1\}$ and $\mathbb{T}=\partial \mathbb{D}$. 
When given a continuous function $f$ on $\mathbb{T}$, the famous \textbf{Dirichlet
problem}, which is to find the solution for the following equations:
\begin{equation*}
  \begin{cases}
  \Delta u=0, & z\in \mathbb{D}\\
  u=f, & z\in \mathbb{T}.
  \end{cases}
\end{equation*}
is asking for a harmonic extension $u$ over $\mathbb{D}$ with the boundary data $f$ on $\mathbb{T}$.
To give a solution for this problem, we'd like to introduce the notion called \textbf{Poisson integral}.
This section is based on \cite{garnett2006bounded}, I.3.
\begin{definition}
Let $c\in\mathbb{C}$. The \textbf{Herglotz--Riesz kernel centered at $c$} is the function
\begin{equation*}
  K(c;z,w) := \frac{(w-c)+(z-c)}{(w-c)-(z-c)} = \frac{w+z-2c}{w-z},
  \qquad z,w\in\mathbb{C},\ w\neq z.
\end{equation*}
The \textbf{Poisson kernel centered at $c$} is defined as
\begin{equation*}
  P(c;z,w) := \text{Re}\; K(c;z,w)= \frac{|w-c|^2-|z-c|^2}{|w-z|^2},
  \qquad z,w\in\mathbb{C},\ w\neq z.
\end{equation*}
and we denote $K(z,w):= K(0;z,w)$, $P(z,w):= P(0;z,w)$, and $K_z(\theta):= K(z,e^{i\theta})$,
 $P_z(\theta):= P(z,e^{i\theta})$. 
\end{definition}

\begin{definition}
Given a function $f(\theta)\in L^1(\mathbb{T})$, we define the \textbf{Poisson integral} of $f$ 
to be the function
\begin{equation*}
  P[f](z):=\frac{1}{2\pi}\int P_z(\theta) f(\theta)d\theta.
\end{equation*}
In general, given a finite (signed) measure $\mu$ on $\mathbb{T}$, we define the \textbf{Poisson integral} of $\mu$
to be the function
\begin{equation*}
  P[\mu](z):=\int P_z(\theta) d\mu (\theta).
\end{equation*}
The first definition is a special case of the second one if we set 
$d\mu(\theta)=\frac{1}{2\pi} f(\theta)d\theta$.
\end{definition}

\begin{theorem}[Poisson integral formula]\label{Poisson integral formula}
Let $u$ be a continuous function on the closed unit disc $\overline{\mathbb{D}}$. If $u$ is
harmonic on $\mathbb{D}$, then
\begin{equation*}
  u(z)=\frac{1}{2\pi}\int_0^{2\pi} P_z(\theta)\, u(e^{i\theta})\, d\theta
  \qquad \text{for all } z\in\mathbb{D}.
\end{equation*}
\end{theorem}
\begin{proof}
Recall that if $u$ is harmonic on $\mathbb{D}$, then it has the 
\textbf{mean value property}: 
\begin{equation*}
  u(0)=\frac{1}{2\pi}\int_0^{2\pi} u(e^{i\theta})d\theta
\end{equation*}
Now given $z_0=re^{i\theta_0}$, consider the M\"{o}bius transformation 
$\tau(z)=(z-z_0)/(1-\bar{z}_0 z)$, and write $\tau(e^{i\theta})=e^{i\varphi}$.
Then we can compute that $d\varphi/d\theta=P_{z_0}(\theta)$, therefore
\begin{equation*}
  u(z_0)=u(\tau^{-1}(0))=\frac{1}{2\pi}\int_0^{2\pi} (u\circ\tau^{-1})(e^{i\varphi})
  d\varphi=\frac{1}{2\pi}\int_0^{2\pi} u(e^{i\theta})P_{z_0}(\theta)d\theta.
\end{equation*}
\end{proof}
Moreover, we have part of the converse of above theorem:
\begin{theorem}\label{Poisson integral of L^1 is harmonic}
Given a function $f(\theta)\in L^1(\mathbb{T})$, its Poisson integral $P[f]$
is harmonic on $\mathbb{D}$.
\end{theorem}
\begin{proof}[Sketch of Proof]
Recall that a function $u(z)$ is said to be harmonic if 
\begin{equation*}
  \Delta u=\frac{\partial^2 u}{\partial x^2}+\frac{\partial^2 u}{\partial y^2}=0.
\end{equation*}
As Herglotz--Riesz kernel $K_z(\theta)$ is analytic over $\mathbb{D}$, 
the Poisson kernel $P_z(\theta)=\text{Re} K_z(\theta)$ is harmonic over 
$\mathbb{D}$, thus for $z\in \mathbb{D}$,
\begin{equation*}
  \Delta P[f](z)=\frac{1}{2\pi}\int \cancel{\Delta P_z(\theta)} f(\theta)d\theta=0.
\end{equation*}
\end{proof}
More generally, the same holds for Poisson integrals of measures:
\begin{theorem}
Given a finite measure $\mu$ on $\mathbb{T}$, its Poisson integral $P[\mu]$
is harmonic on $\mathbb{D}$.
\end{theorem}
\begin{proof}[Sketch of Proof]
Similar as in the previous proof, for $z\in\mathbb{D}$,
\begin{equation*}
  \Delta P[\mu](z)=\int \cancel{\Delta P_z(\theta)}\, d\mu(\theta)=0.
\end{equation*}
\end{proof}
The full converse theorem of Poisson integral formula states that the Poisson integral indeed
gives a continuous solution for Dirichlet problem, which will be later proved in Corollary 
\ref{Dirichlet problem on the disc}. Before that, we'd like to introduce more properties of 
Poisson integral.



Here are some properties for Poisson kernel:
\begin{lemma}\label{Properties of Poisson kernel}
Let $z\in\mathbb{D}$. Then
\begin{itemize}
  \item[(a)] $P_z(\theta)>0$ for all $\theta$, and
  $\displaystyle\frac{1}{2\pi}\int_0^{2\pi} P_z(\theta)\, d\theta=1$;
  \item[(b)] for any $\theta_0\in\mathbb{R}$ and $0<\delta<\pi$,
  \begin{equation*}
    \sup_{\delta\leq|\theta-\theta_0|\leq\pi}P_z(\theta)\rightarrow 0
    \quad\text{and}\quad
    \frac{1}{2\pi}\int_{\delta\leq|\theta-\theta_0|\leq\pi}P_z(\theta)\, d\theta\rightarrow 0
    \qquad (z\rightarrow e^{i\theta_0}).
  \end{equation*}
\end{itemize}
\end{lemma}

Here's an important estimate for Poisson integral of $L^p$ function:
\begin{theorem}
Let $1\leq p\leq\infty$ and let $f$ be a function in $L^p(\mathbb{T})$, and
$\|f\|_p^p:=\frac{1}{2\pi}\int_0^{2\pi}|f(\theta)|^p\, d\theta$ for $p<\infty$. Then for every
$0<r<1$,
\begin{equation*}
  \|P[f](re^{i\cdot})\|_p\leq\|f\|_p,
\end{equation*}
and hence $\sup_{0<r<1}\|P[f](re^{i\cdot})\|_p\leq\|f\|_p$.
\end{theorem}
\begin{proof}[Sketch of Proof]
If $p=\infty$, then $|P[f](z)|\leq\frac{1}{2\pi}\int P_z(\theta)|f(\theta)|\, d\theta
\leq\|f\|_\infty$.

If $1\leq p<\infty$, for $z\in\mathbb{D}$, part (a) of the lemma \ref{Properties of Poisson kernel} says that $P_z(\theta)\frac{d\theta}{2\pi}$ is a
probability measure on $\mathbb{T}$. Jensen's inequality for this probability measure gives
\begin{equation*}
  |P[f](z)|^p\leq\frac{1}{2\pi}\int_0^{2\pi} P_z(\theta)|f(\theta)|^p\, d\theta.
\end{equation*}
Set $z=re^{i\varphi}$. Note the symmetry of Poisson kernel:
\begin{equation*}
  P_{re^{i\varphi}}(\theta)=\frac{1-r^2}{1-2r\cos(\theta-\varphi)+r^2}=P_{re^{i\theta}}(\varphi),
\end{equation*}
Thus:
\begin{equation*}
  \frac{1}{2\pi}\int_0^{2\pi}|P[f](re^{i\varphi})|^p\, d\varphi
  \leq\frac{1}{2\pi}\int_0^{2\pi}|f(\theta)|^p
  \left(\frac{1}{2\pi}\int_0^{2\pi}P_{re^{i\theta}}(\varphi)\, d\varphi\right) d\theta
  =\|f\|_p^p. \qedhere
\end{equation*}
\end{proof}

The following theorem shows that the Poisson kernel acts like the identity as $r\rightarrow 1$:
\begin{theorem}\label{Poisson approximate identity}
Let $C(\mathbb{T})$ be the space of all continuous functions over $\mathbb{T}$. Then
\begin{itemize}
  \item[(a)] If $f\in C(\mathbb{T})$, then $P[f](re^{i\theta})$ converges uniformly to $f(e^{i\theta})$
  as $r\rightarrow 1$.
  \item[(b)] If $\mu$ is a finite measure on $\mathbb{T}$, then the measures
  $P[\mu](re^{i\theta})\, d\theta/(2\pi)$ converge weak-star to $\mu$ as $r\rightarrow 1$, i.e.
  \begin{equation*}
    \frac{1}{2\pi}\int_0^{2\pi} g(\theta)P[\mu](re^{i\theta})\, d\theta\rightarrow\int g\, d\mu
    \qquad (r\rightarrow 1)
  \end{equation*}
  for every $g\in C(\mathbb{T})$.
\end{itemize}
\end{theorem}
\begin{proof}[Sketch of Proof]
(a) Fix $\varphi$. By part (a) of the lemma \ref{Properties of Poisson kernel},
\begin{equation*}
  |P[f](re^{i\varphi})-f(e^{i\varphi})|
  \leq\frac{1}{2\pi}\int_{|\theta-\varphi|<\delta}P_{re^{i\varphi}}(\theta)|f(e^{i\theta})-f(e^{i\varphi})|\, d\theta
  +\frac{1}{2\pi}\int_{\delta\leq|\theta-\varphi|\leq\pi}P_{re^{i\varphi}}(\theta)|f(e^{i\theta})-f(e^{i\varphi})|\, d\theta.
\end{equation*}
As $\mathbb{T}$ is compact, $f$ is uniformly continuous on $\mathbb{T}$.
Given $\varepsilon>0$, by uniform continuity of $f$, we can pick $\delta>0$ s.t.
$|f(e^{i\theta})-f(e^{i\varphi})|<\varepsilon$ for all $\varphi$ and all $|\theta-\varphi|<\delta$,
then we have
\begin{equation*}
  \frac{1}{2\pi}\int_{|\theta-\varphi|<\delta}P_{re^{i\varphi}}(\theta)|f(e^{i\theta})-f(e^{i\varphi})|\, d\theta
  \leq \varepsilon\cdot\frac{1}{2\pi}\int P_{re^{i\varphi}}(\theta)d\theta=\varepsilon.
\end{equation*}
For this $\delta$, the
second term is at most $2\|f\|_\infty\cdot\frac{1}{2\pi}\int_{\delta\leq|\theta-\varphi|\leq\pi}
P_{re^{i\varphi}}(\theta)\, d\theta$. Since $P_{re^{i\varphi}}(\theta)=P_r(\theta-\varphi)$, this integral
does not depend on $\varphi$, and it tends to $0$ as $r\rightarrow 1$ by part (b) of the
lemma \ref{Properties of Poisson kernel}. Hence the convergence is uniform in $\varphi$.

(b) Let $g\in C(\mathbb{T})$. For fixed $r$, the Poisson kernel is bounded, so by Fubini's theorem and
the symmetry $P_{re^{i\theta}}(\varphi)=P_{re^{i\varphi}}(\theta)$,
\begin{equation*}
  \frac{1}{2\pi}\int_0^{2\pi} g(\theta)P[\mu](re^{i\theta})\, d\theta
  =\int\left(\frac{1}{2\pi}\int_0^{2\pi}P_{re^{i\varphi}}(\theta)g(\theta)\, d\theta\right) d\mu(\varphi)
  =\int P[g](re^{i\varphi})\, d\mu(\varphi).
\end{equation*}
By (a), $P[g](re^{i\varphi})\rightarrow g(e^{i\varphi})$ uniformly, and $\mu$ is finite, so the
right-hand side tends to $\int g\, d\mu$.
\end{proof}

As a consequence, the Poisson integral solves the Dirichlet problem on the disc (which is 
the converse result of the Poisson integral formula
 \ref{Poisson integral formula}):
\begin{corollary}[Dirichlet problem on the disc]\label{Dirichlet problem on the disc}
Let $f$ be a continuous function on $\mathbb{T}$. Define
\begin{equation*}
  u(z):=
  \begin{cases}
    P[f](z), & z\in\mathbb{D},\\
    f(z), & z\in\mathbb{T}.
  \end{cases}
\end{equation*}
Then $u$ is continuous on $\overline{\mathbb{D}}$ and harmonic on $\mathbb{D}$, i.e. 
a continuous solution on $\overline{\mathbb{D}}$ for the Dirichlet
problem:
\begin{equation*}
  \begin{cases}
  \Delta u=0, & z\in \mathbb{D}\\
  u=f, & z\in \mathbb{T}.
  \end{cases}
\end{equation*}
\end{corollary}
\begin{proof}
$u$ is harmonic on $\mathbb{D}$ by Theorem \ref{Poisson integral of L^1 is harmonic}. For
continuity at $e^{i\theta_0}\in\mathbb{T}$, let $z=re^{i\varphi}\in\mathbb{D}$; then
\begin{equation*}
  |u(z)-f(e^{i\theta_0})|\leq|P[f](re^{i\varphi})-f(e^{i\varphi})|+|f(e^{i\varphi})-f(e^{i\theta_0})|,
\end{equation*}
which tends to $0$ as $z\rightarrow e^{i\theta_0}$ by part (a) of
Theorem \ref{Poisson approximate identity} and the continuity of $f$.
\end{proof}

Moreover, a measure on $\mathbb{T}$ is determined by its Poisson integral:
\begin{corollary}
Let $\mu$ and $\nu$ be finite measures on $\mathbb{T}$. If $P[\mu]=P[\nu]$ on $\mathbb{D}$,
then $\mu=\nu$.
\end{corollary}
\begin{proof}
Let $\lambda:=\mu-\nu$. By linearity, $P[\lambda]=0$ on $\mathbb{D}$. For every $g\in C(\mathbb{T})$,
part (b) of Theorem \ref{Poisson approximate identity} gives
\begin{equation*}
  \int g\, d\lambda=\lim_{r\rightarrow 1}\frac{1}{2\pi}\int_0^{2\pi} g(\theta)P[\lambda](re^{i\theta})\, d\theta=0.
\end{equation*}
By the uniqueness part of the Riesz representation theorem, $\lambda=0$, thus $\mu=\nu$.
\end{proof}

Finally, Poisson integrals of $L^p$ functions and of measures are characterized by the growth
of their radial means:
\begin{theorem}\label{Characterization of Poisson integrals}
Let $u$ be a harmonic function on $\mathbb{D}$.
\begin{itemize}
  \item[(a)] If $1<p\leq\infty$, then $u=P[f]$ for some $f\in L^p(\mathbb{T})$ if and only if
  \begin{equation*}
    \sup_{0<r<1}\|u(re^{i\cdot})\|_p<\infty.
  \end{equation*}
  \item[(b)] $u=P[\mu]$ for some finite measure $\mu$ on $\mathbb{T}$ if and only if
  \begin{equation*}
    \sup_{0<r<1}\|u(re^{i\cdot})\|_1<\infty.
  \end{equation*}
  \item[(c)] (Herglotz) $u\geq 0$ on $\mathbb{D}$ if and only if $u=P[\mu]$ for some finite
  positive measure $\mu$ on $\mathbb{T}$.
\end{itemize}
\end{theorem}

\begin{proof}
(\textcolor{red}{missing})
\end{proof}

\section{Nontangential limit}
When we are discussing functions over $\mathbb{D}$, it will be good to have the proper
notion of `limit function' over the boundary $\mathbb{T}=\partial \mathbb{D}$.
This section is mainly based on \cite{rudin1974real}, 11.21 - 11.24.
\begin{definition}
Fix $\alpha>1$. For $\zeta\in\mathbb{T}$, we consider the \textbf{Stolz region}:
\begin{equation*}
  \Gamma_\alpha(\zeta):=\left\{z\in\mathbb{D}: |\zeta-z|<\alpha(1-|z|)\right\}.
\end{equation*}
A function $u:\mathbb{D}\to\mathbb{C}$ is said to have \textbf{nontangential limit} $\lambda\in\mathbb{C}$
at $\zeta\in\mathbb{T}$ if, for every $\alpha>1$,
\begin{equation*}
  \lim_{\substack{z\to\zeta\\ z\in\Gamma_\alpha(\zeta)}}u(z)=L.
\end{equation*}
When this limit exists, we write $u^*(\zeta)=\lambda$ and call it the \textbf{(nontangential) boundary value} of $u$ at $\zeta$.
This give rise to a function $u^*:\mathbb{T}\rightarrow\mathbb{C}$, for which we call the \textbf{(nontangential) boundary function}
of $u$.
\end{definition}


Before we discuss the main results related with nontangential limit, we first 
introduce a notion called \textbf{Lebesgue point}:
\begin{definition}[Lebesgue point]
Consider $(\mathbb{T},d\sigma)$, where $d\sigma := d\theta/(2\pi)$.
If $f\in L^1(\mathbb{T})$, a point $\zeta=e^{i\theta_0}\in\mathbb{T}$
is called a \textbf{Lebesgue point} of $f$ if
\begin{equation*}
  \lim_{h\downarrow 0}
  \frac{1}{2h}
  \int_{\theta_0-h}^{\theta_0+h}
  \left|f(e^{i\theta})-f(e^{i\theta_0})\right|\,d\theta=0.
\end{equation*}
\end{definition}

\begin{remark}
If $f$ is continuous at $\xi$, then $\xi$ is a Lebesgue point of $f$. So one can
think of Lebesgue points of a function are those points where in average sense, the
function does not oscillate too much.
\end{remark}

\begin{theorem}[Lebesgue differentiation theorem]
\label{thm:lebesgue_differentiation}
Let $f\in L^1(\mathbb{T})$. Then almost every point of $\mathbb{T}$
 is a Lebesgue point of $f$.
\end{theorem}

The following result tells us that nontangential limit of
the Poisson integral of a function gives back the function value itself
at all lebesgue points of that function:
\begin{theorem}
\label{thm:poisson_nontangential_lebesgue_point}
Let $f\in L^1(\mathbb{T})$ be real-valued. If $\zeta\in\mathbb{T}$ is a
Lebesgue point of $f$, then the Poisson integral $P[f]$ has nontangential
limit $f(\zeta)$ at $\zeta$. That is, for every $\alpha>1$,
\begin{equation*}
  \lim_{\substack{z\to\zeta\\ z\in\Gamma_\alpha(\zeta)}}P[f](z)=f(\zeta).
\end{equation*}
\end{theorem}

Then using Lebesgue differentiation theorem \ref{thm:lebesgue_differentiation},
we get an a.e. recovery of boundary data:
\begin{corollary}[Almost everywhere recovery of boundary data]
\label{cor:poisson_nontangential_ae}
Let $f\in L^1(\mathbb{T})$ be real-valued. Then $P[f]$ has nontangential
limits almost everywhere on $\mathbb{T}$, and
\begin{equation*}
  (P[f])^*(\zeta)=f(\zeta)
  \qquad\text{for almost every }\zeta\in\mathbb{T}.
\end{equation*}
\end{corollary}

\begin{theorem}
\label{thm:singular_poisson_nontangential_zero}
Let $\nu$ be a finite measure on $\mathbb{T}$ s.t. $\nu\bot d\sigma$, i.e. singular
with respect to $d\sigma = d\theta/(2\pi)$. Then its Poisson integral $P[\nu]$ has
nontangential limit 0 at almost every point of $\mathbb{T}$. For every $\alpha>1$,
\begin{equation*}
  \lim_{\substack{z\to\zeta\\ z\in\Gamma_\alpha(\zeta)}}P[\nu](z)=0
  \qquad\text{for a.e. }\zeta\in\mathbb{T}.
\end{equation*}
\end{theorem}

\begin{theorem}[Fatou's theorem for finite measures]
\label{thm:fatou_signed_measure}
Let $\mu$ be a finite measure on $\mathbb{T}$ with Lebesgue
decomposition
\begin{equation*}
  d\mu=f\,\frac{d\theta}{2\pi}+d\mu_s,
  \qquad f\in L^1(\mathbb{T}),
  \qquad \mu_s\perp\frac{d\theta}{2\pi}.
\end{equation*}
Then $P[\mu]$ has nontangential limits a.e. on $\mathbb{T}$, and for every $\alpha>1$,
\begin{equation*}
  \lim_{\substack{z\to\zeta\\ z\in\Gamma_\alpha(\zeta)}}P[\mu](z)=f(\zeta)
  \qquad\text{for a.e. }\zeta\in\mathbb{T}.
\end{equation*}
\end{theorem}
\begin{proof}
By linearity, $P[\mu]=P[f]+P[\mu_s]$. Apply
Corollary~\ref{cor:poisson_nontangential_ae} and
Theorem~\ref{thm:singular_poisson_nontangential_zero} will get the desired result.
\end{proof}





\section{Subharmonic function}
\label{subsec:subharmonic}
To investigate the main objects in this thesis,
which are Hardy spaces and Nevanlinna class, it will be good to know the properties of 
function like $|f|^p$, $\log |f|$ and $\log^+ |f|$
for analytic function $f$. These functions belong to an important class of functions called 
\textbf{subharmonic functions}. This section is based on \cite{garnett2006bounded}, I.6.
\begin{definition}
Let $\Omega$ be an open set in the plane. A \textbf{subharmonic function} on $\Omega$ is a function
$v : \Omega \to [-\infty, \infty)$ such that
\begin{itemize}
  \item[(a)] $v$ is \emph{upper semicontinous}, i.e. $v(z_0)\geq \lim_{z\rightarrow z_0}v(z)$ 
  for all $z_0\in\Omega$;
  \item[(b)] for each $z_0 \in \Omega$ there is $r(z_0) > 0$ such that the disc 
$B(z_0, r(z_0)) = \{z : |z - z_0| < r(z_0)\}$ is contained in $\Omega$ and 
such that for every $r < r(z_0)$,
\begin{equation}\label{MVP of subharmonic functions}
  v(z_0) \leq \frac{1}{\pi r^2} \iint_{B(z_0,r)} v(z)\, dx\, dy.
\end{equation}
We often call (\ref{MVP of subharmonic functions}) the `\textbf{submean value property}'.
\end{itemize}
\end{definition}
\begin{remark}
Note that the integral in \ref{MVP of subharmonic functions} either
converges or diverges to $-\infty$, due to the upper-semicontinuity by the following
lemma:
\begin{lemma}\label{boundedness of uppersemicont fct}
Let $X$ be a topological space.
If $v:X\rightarrow [-\infty,\infty)$ is upper-semicontinuous, then $v$ is bounded
above on each compact subset $K\subset X$.
\end{lemma}
\end{remark}

Here are some examples of subharmonic functions:
\begin{example}
Every harmonic function is subharmonic, due to the \textbf{mean value property} of harmonic function:
\begin{theorem}
Let $\Omega$ be an open set in the plane and let $u$ be a harmonic function on $\Omega$. Then for
every $z_0\in\Omega$ and every $r>0$ such that the closed disc $\overline{B(z_0,r)}$ is contained
in $\Omega$,
\begin{equation*}
  u(z_0) = \frac{1}{2\pi}\int_0^{2\pi} u(z_0+re^{i\theta})\, d\theta
  = \frac{1}{\pi r^2} \iint_{B(z_0,r)} u(z)\, dx\, dy.
\end{equation*}
\end{theorem}
\end{example}
\begin{example}\label{log |f| is subharmonic}
Let $f(z)$ be an analytic function on a domain $\Omega\subset \mathbb{C}$, then 
$v(z)=\log |f(z)|$ is subharmonic on $\Omega$: 
\begin{itemize}
  \item $v$ is upper semicontinuous;
  \item For $z_0$ s.t. $f(z_0)=0$, condition \ref{MVP of subharmonic functions} trivially
  satisfied;
  \item For $z_0$ s.t. $f(z_0)\neq 0$, we can write $v(z)=\text{Re}(\log f(z))$ on some neighborhood
  of $z_0$, which is harmonic on this neighborhood.
\end{itemize}
\end{example}
\begin{example}
Let $f(z)$ be an analytic function on a domain $\Omega\subset \mathbb{C}$, then 
$|f(z)|^p$ is subharmonic on $\Omega$ as 
\begin{equation*}
  |f(z)|^p=\exp(p \log |f(z)|)
\end{equation*}
So the condition \ref{MVP of subharmonic functions} can be proved by the fact that
$\log |f(z)|$ is subharmonic and Jensen's inequality as $\exp$ is convex. More generally, we
have the following lemma:
\begin{lemma}
Given a subharmonic funciton $v$ on $\Omega$ and an increasing convex function $\varphi$ on
$[-\infty, \infty)$ which is continuous at $-\infty$, then $\varphi\circ v$ is subharmonic on
$\Omega$.
\end{lemma}
With this lemma, we can see that
\begin{equation*}
    \log^+ |f(z)|=\max(\log |f(z)|,0)
\end{equation*}
is also a subharmonic function.
\end{example}

There is an equivalent characterization of subharmonicity:
\begin{theorem}[Harmonic comparison characterization]\label{Harmonic comparison characterization}
Let $u:\Omega\rightarrow [-\infty,\infty)$ be an upper semicontinuous function. Then $u$ is 
subharmonic on $\Omega$ if and only if the following condition holds: If $v$ is a continuous function
on $\overline{B(z_0,r)}\subset \Omega$ that is harmonic on $B(z_0,r)$ s.t. for all 
$\xi\in \partial B(z_0,r)$, we have $u(\xi)\leq v(\xi)$, then 
\begin{equation*}
  u(z)\leq v(z),\quad z\in B(z_0,r).
\end{equation*}
\end{theorem}
\begin{proof}
(\textcolor{red}{missing})
\end{proof}

\begin{corollary}
If $\Omega$ is a connected open set and $v$ is a subharmonic function on $\Omega$ s.t. $v\not\equiv -\infty$,
then given any closed disc $\overline{B(z_0,r)}\subset \Omega$, 
\begin{equation*}
    \frac{1}{2\pi}\int v(z_0+r e^{i\theta})d\theta \neq -\infty.
\end{equation*}
\end{corollary}
\begin{proof}[Sketch of Proof]
First recall a lemma:
\begin{lemma}
Let $X$ be a metric space, $v$ be an upper-semicontinuous function. Then we can always find
a sequence $u_n$ of continuous function that decreasing to $v$ on $X$.
\end{lemma}
Now we pick $u_n$ to be a sequence of continuous functions decreasing to $v$ on 
$\partial B(z_0,r)$. Let $U_n(z)=P[u_n](z)$, i.e. the harmonic extension of $u_n$ to $B(z_0,r)$.
Assume for the sake of contradiction that we can find some $\overline{B(z_0,r)}\subset\Omega$ s.t.
\begin{equation*}
  \frac{1}{2\pi}\int v(z_0+r e^{i\theta})d\theta = -\infty.
\end{equation*}
Recall Lemma \ref{boundedness of uppersemicont fct}, $v$ is bounded above by some
constant $M\geq 0$. Write $v=v^+-v^-$ with $v^+=\max(v,0)$ and $v^-=\max(-v,0)$, so that
$0\leq v^+\leq M$ and $0\leq v^-\leq\infty$ on $\partial B(z_0,r)$. Since
\begin{equation*}
  \frac{1}{2\pi}\int v^+(z_0+re^{i\theta})\, d\theta\leq M<\infty,
\end{equation*}
the assumption is equivalent to
\begin{equation*}
  \frac{1}{2\pi}\int v^-(z_0+re^{i\theta})\, d\theta=\infty.
\end{equation*}
Now fix $z$ with $|z|<1$. Since $1-|z|\leq |e^{i\theta}-z|\leq 1+|z|$, we have
\begin{equation*}
  0<\frac{1-|z|}{1+|z|}\leq P_z(\theta)=\frac{1-|z|^2}{|e^{i\theta}-z|^2}\leq\frac{1+|z|}{1-|z|}
  \quad\text{for all }\theta.
\end{equation*}
Therefore
\begin{align*}
  \frac{1}{2\pi}\int P_z(\theta)\, v^+(z_0+re^{i\theta})\, d\theta
  &\leq \frac{1+|z|}{1-|z|}\, M<\infty,\\
  \frac{1}{2\pi}\int P_z(\theta)\, v^-(z_0+re^{i\theta})\, d\theta
  &\geq \frac{1-|z|}{1+|z|}\cdot\frac{1}{2\pi}\int v^-(z_0+re^{i\theta})\, d\theta=\infty.
\end{align*}
Therefore,
\begin{equation*}
  \frac{1}{2\pi}\int P_z(\theta)\, v(z_0+re^{i\theta})\, d\theta = -\infty,\quad |z|<1.
\end{equation*}
As
\begin{equation*}
  U_n(z)=\frac{1}{2\pi}\int P_z(\theta)\, u_n(z_0+re^{i\theta})\, d\theta 
\end{equation*}
and $u_n\downarrow v$, so $U_n(z)\rightarrow -\infty$ for each $z\in B(z_0,r)$. By Harmonic
comparison \ref{Harmonic comparison characterization}, we know that $U_n \geq v$ over $B(z_0,r)$,
so $v=-\infty$ on $B(z_0,r)$. With above result, we can actually extend this local property globally
using standard topological argument as follow: consider 
\begin{equation*}
  A:=\{z\in \Omega: v(z)\equiv -\infty\text{ on a neighborhood of }z\}
\end{equation*}
then one can show that it is open and closed. As $z_0\in A$, it is nonempty, thus $v\equiv -\infty$
over $\Omega$, which gives a contradiction.

\end{proof}

To investigate further properties of subharmonic functions, we'd like to first define
an auxiliary function called \textbf{Poisson modification}:
\begin{definition}[Poisson modification]\label{Poisson modification}
Let $v:\mathbb{D}\rightarrow[-\infty,\infty)$ be an upper semicontinuous function and let
$0<r<1$. The \textbf{Poisson modification} of $v$ on $B(0,r)$ is the function
\begin{equation*}
  v_r(z):=
  \begin{cases}
    v(z), & r\leq|z|<1,\\[4pt]
    \displaystyle\frac{1}{2\pi}\int P_{z/r}(\theta)\, v(re^{i\theta})\, d\theta, & |z|<r.
  \end{cases}
\end{equation*}
\end{definition}
By Lemma \ref{boundedness of uppersemicont fct}, $v$ is bounded above on $\partial B(0,r)$, so
the integral is well-defined with values in $[-\infty,\infty)$.

\begin{theorem}\label{Properties of Poisson modification}
Let $v$ be a subharmonic function on $\mathbb{D}$ s.t. $v\not\equiv-\infty$, and let $0<r<1$.
Then
\begin{itemize}
  \item[(a)] $v_r$ is finite and harmonic on $B(0,r)$;
  \item[(b)] $v_r$ is subharmonic on $\mathbb{D}$;
  \item[(c)] $v(z)\leq v_r(z)$ for all $z\in\mathbb{D}$;
  \item[(d)] $v_r(z)$ is an increasing function of $r$.
\end{itemize}
\end{theorem}
\begin{proof}
(\textcolor{red}{missing})
\end{proof}
If we take $z=0$ in Theorem \ref{Properties of Poisson modification}(d), we get:
\begin{corollary}
If $v(z)$ is a subharmonic function on $\mathbb{D}$, then
\begin{equation*}
  m(r)=v_r(0)=\frac{1}{2\pi}\int v(r e^{i\theta})d\theta
\end{equation*}
is an increasing function of $r\in(0,1)$.
\end{corollary}

\begin{definition}[Harmonic majorant]\label{Harmonic majorant}
Let $v$ be a subharmonic function on an open set $\Omega$. A \textbf{harmonic majorant} of $v$
is a harmonic function $U$ on $\Omega$ s.t. $v(z)\leq U(z)$ for all $z\in\Omega$. A harmonic
majorant $u$ of $v$ is called the \textbf{least harmonic majorant} of $v$ if $u(z)\leq U(z)$,
$z\in\Omega$, for every harmonic majorant $U$ of $v$.
\end{definition}

\begin{theorem}\label{Harmonic majorant criterion}
Let $v$ be a subharmonic function on $\mathbb{D}$ s.t. $v\not\equiv-\infty$. Then $v$ has a
harmonic majorant if and only if
\begin{equation*}
  \sup_{0<r<1}\frac{1}{2\pi}\int v(re^{i\theta})\, d\theta=\sup_{0<r<1}v_r(0)<\infty.
\end{equation*}
In this case, the least harmonic majorant of $v$ is
\begin{equation*}
  u(z)=\lim_{r\rightarrow 1}\frac{1}{2\pi}\int P_{z/r}(\theta)\, v(re^{i\theta})\, d\theta
  =\lim_{r\rightarrow 1}v_r(z),\quad z\in\mathbb{D}.
\end{equation*}
\end{theorem}

\begin{proof}
\textcolor{red}{missing}
\end{proof}